% 2
\vspace{0.4em}

\begin{center}
\textbf{Constraint (2)}
\end{center}

\noindent
First constituents that are constituted by simplex masculine or neuter nouns that build the plural form with a zero ending and simplex or complex feminine nouns that build the plural form with a zero ending attach a zero linker almost regularly.
\vspace{0.5em}

\begin{tikzpicture}
\node[
    draw,
    rounded corners=6pt,
    fill=gray!30,
    inner sep=2pt
] (11) {
    \colorbox{gray!30}{P2L}
};
\node[
    draw,
    rounded corners=6pt,
    fill=gray!30,
    inner sep=2pt,
    right=0.5em of 11
] (12) {
    \colorbox{gray!30}{MD}
};
\node[
    draw,
    rounded corners=6pt,
    fill=gray!30,
    inner sep=2pt,
    right=0.5em of 12
] {
    \colorbox{gray!30}{first constituent}
};
\end{tikzpicture}

\vspace{-0.25em}

\begin{table}[!ht]
\begin{center}
\setlength{\extrarowheight}{1pt}
\begin{tabularx}{\columnwidth}{|X|X|X|}
\hline
 & item & type \\
\hline
 coverage & $0.067$ & $0.066$ \\
\hline
 regularity & $0.954$ & $0.936$ \\
\hline
\end{tabularx}
\end{center}
\end{table}

\vspace{-1.35em}

\noindent
Support level: \textit{supported}

\vspace{0.3em}

\noindent
Examples: Schlitten{\zeroLinker}fahrt, Wagen{\zeroLinker}rad, Ufer{\zeroLinker}promenade

\noindent
Counterexamples: Orden\textbf{s}bruder, Teufel\textbf{s}werk, Frieden\textbf{s}angebot

\vspace{0.3em}

\noindent
Sources: \citet[106]{ortner_mueller-bollhagen_1991}

\vspace{0.35em}
\hrule




% 3
\vspace{0.4em}

\begin{center}
\textbf{Constraint (3)}
\end{center}

\noindent
First constituents that are constituted by nouns that build the plural form with -s attach a zero linker regularly.
\vspace{0.5em}

\begin{tikzpicture}
\node[
    draw,
    rounded corners=6pt,
    fill=gray!30,
    inner sep=2pt
] (11) {
    \colorbox{gray!30}{P2L}
};
\node[
    draw,
    rounded corners=6pt,
    fill=gray!30,
    inner sep=2pt,
    right=0.5em of 11
] (12) {
    \colorbox{gray!30}{M}
};
\node[
    draw,
    rounded corners=6pt,
    fill=gray!30,
    inner sep=2pt,
    right=0.5em of 12
] {
    \colorbox{gray!30}{first constituent}
};
\end{tikzpicture}

\vspace{-0.25em}

\begin{table}[!ht]
\begin{center}
\setlength{\extrarowheight}{1pt}
\begin{tabularx}{\columnwidth}{|X|X|X|}
\hline
 & item & type \\
\hline
 coverage & $0.035$ & $0.033$ \\
\hline
 regularity & $0.97$ & $0.957$ \\
\hline
\end{tabularx}
\end{center}
\end{table}

\vspace{-1.35em}

\noindent
Support level: \textit{partially supported}

\vspace{0.3em}

\noindent
Examples: Balkon{\zeroLinker}möbel, Hotel{\zeroLinker}kette, Bob{\zeroLinker}bahn

\noindent
Counterexamples: Uhu\textbf{s}nest, Deck\textbf{s}balken, Interim\textbf{s}phase, Karneval\textbf{s}kostüm

\vspace{0.3em}

\noindent
Sources: \citet[227]{eisenberg_2013}, \citet[572]{fuhrhop_kuerschner_2015}, \citet[9]{kuerschner_2010}, \citet[78]{nuebling_szczepaniak_2013}, \citet[83, 106]{ortner_mueller-bollhagen_1991}, \citet[229]{schaefer_2018}

\vspace{0.35em}
\hrule




% 4
\vspace{0.4em}

\begin{center}
\textbf{Constraint (4)}
\end{center}

\noindent
First constituents that are constituted by nouns that build the plural form with -e mostly have a zero linker.
\vspace{0.5em}

\begin{tikzpicture}
\node[
    draw,
    rounded corners=6pt,
    fill=gray!30,
    inner sep=2pt
] (11) {
    \colorbox{gray!30}{P2L}
};
\node[
    draw,
    rounded corners=6pt,
    fill=gray!30,
    inner sep=2pt,
    right=0.5em of 11
] (12) {
    \colorbox{gray!30}{M}
};
\node[
    draw,
    rounded corners=6pt,
    fill=gray!30,
    inner sep=2pt,
    right=0.5em of 12
] {
    \colorbox{gray!30}{first constituent}
};
\end{tikzpicture}

\vspace{-0.25em}

\begin{table}[!ht]
\begin{center}
\setlength{\extrarowheight}{1pt}
\begin{tabularx}{\columnwidth}{|X|X|X|}
\hline
 & item & type \\
\hline
 coverage & $0.159$ & $0.184$ \\
\hline
 regularity & $0.762$ & $0.689$ \\
\hline
\end{tabularx}
\end{center}
\end{table}

\vspace{-1.35em}

\noindent
Support level: \textit{supported}

\vspace{0.3em}

\noindent
Examples: Tisch{\zeroLinker}decke, Projekt{\zeroLinker}gruppe, Pilz{\zeroLinker}sammler

\noindent
Counterexamples: Tag\textbf{e}gericht, Projekt\textbf{e}macherei, Hund\textbf{e}leine

\vspace{0.3em}

\noindent
Sources: \citet[83]{ortner_mueller-bollhagen_1991}

\vspace{0.35em}
\hrule




% 5
\vspace{0.4em}

\begin{center}
\textbf{Constraint (5)}
\end{center}

\noindent
First constituents that are constituted by nouns that build the plural form with -e attach -e- irregularly in compounds in which the second constituent forces a plural meaning of the given first constituent.
\vspace{0.5em}

\begin{tikzpicture}
\node[
    draw,
    rounded corners=6pt,
    fill=gray!30,
    inner sep=2pt
] (11) {
    \colorbox{gray!30}{P2L}
};
\node[
    draw,
    rounded corners=6pt,
    fill=gray!30,
    inner sep=2pt,
    right=0.5em of 11
] (12) {
    \colorbox{gray!30}{MS}
};
\node[
    draw,
    rounded corners=6pt,
    fill=gray!30,
    inner sep=2pt,
    right=0.5em of 12
] {
    \colorbox{gray!30}{first constituent}
};
\end{tikzpicture}

\vspace{-0.25em}

\begin{table}[!ht]
\begin{center}
\setlength{\extrarowheight}{1pt}
\begin{tabularx}{\columnwidth}{|X|X|X|}
\hline
 & item & type \\
\hline
 coverage & $NA$ & $NA$ \\
\hline
 regularity & $NA$ & $NA$ \\
\hline
\end{tabularx}
\end{center}
\end{table}

\vspace{-1.35em}

\noindent
Support level: \textit{NA}

\vspace{0.3em}

\noindent
Examples: Projekt\textbf{e}macherei, Tag\textbf{e}buch, Punkt\textbf{e}stand

\noindent
Counterexamples: Pilz{\zeroLinker}sammler, Brot{\zeroLinker}korb, Jahr{\zeroLinker}zehnt, Tag\textbf{e}lohn

\vspace{0.3em}

\noindent
Sources: \citet[103]{ortner_mueller-bollhagen_1991}, \citet[29]{schaefer_pankratz_2018}

\vspace{0.35em}
\hrule




% 6
\vspace{0.4em}

\begin{center}
\textbf{Constraint (6)}
\end{center}

\noindent
First constituents that are constituted by monosyllabic animal designations that build the plural form with -e attach -e- irregularly.
\vspace{0.5em}

\begin{tikzpicture}
\node[
    draw,
    rounded corners=6pt,
    fill=gray!30,
    inner sep=2pt
] (11) {
    \colorbox{gray!30}{P2L}
};
\node[
    draw,
    rounded corners=6pt,
    fill=gray!30,
    inner sep=2pt,
    right=0.5em of 11
] (12) {
    \colorbox{gray!30}{MPS}
};
\node[
    draw,
    rounded corners=6pt,
    fill=gray!30,
    inner sep=2pt,
    right=0.5em of 12
] {
    \colorbox{gray!30}{first constituent}
};
\end{tikzpicture}

\vspace{-0.25em}

\begin{table}[!ht]
\begin{center}
\setlength{\extrarowheight}{1pt}
\begin{tabularx}{\columnwidth}{|X|X|X|}
\hline
 & item & type \\
\hline
 coverage & $NA$ & $NA$ \\
\hline
 regularity & $NA$ & $NA$ \\
\hline
\end{tabularx}
\end{center}
\end{table}

\vspace{-1.35em}

\noindent
Support level: \textit{NA}

\vspace{0.3em}

\noindent
Examples: Hund\textbf{e}leine, Swein\textbf{e}stall, Pferd\textbf{e}wagen

\noindent
Counterexamples: Hund\textbf{s}stern, Schwein{\zeroLinker}kram, Schaff{\zeroLinker}stall

\vspace{0.3em}

\noindent
Sources: \citet[228]{eisenberg_2013}, \citet[11]{nuebling_szczepaniak_2008}, \citet[81]{nuebling_szczepaniak_2013}

\vspace{0.35em}
\hrule




% 7
\vspace{0.4em}

\begin{center}
\textbf{Constraint (7)}
\end{center}

\noindent
First constituents that are constituted by nouns that build the plural form with -(")er mostly attach a zero linker or -(")er-.
\vspace{0.5em}

\begin{tikzpicture}
\node[
    draw,
    rounded corners=6pt,
    fill=gray!30,
    inner sep=2pt
] (11) {
    \colorbox{gray!30}{P2L}
};
\node[
    draw,
    rounded corners=6pt,
    fill=gray!30,
    inner sep=2pt,
    right=0.5em of 11
] (12) {
    \colorbox{gray!30}{M}
};
\node[
    draw,
    rounded corners=6pt,
    fill=gray!30,
    inner sep=2pt,
    right=0.5em of 12
] {
    \colorbox{gray!30}{first constituent}
};
\end{tikzpicture}

\vspace{-0.25em}

\begin{table}[!ht]
\begin{center}
\setlength{\extrarowheight}{1pt}
\begin{tabularx}{\columnwidth}{|X|X|X|}
\hline
 & item & type \\
\hline
 coverage & $0.019$ & $0.058$ \\
\hline
 regularity & $0.833$ & $0.868$ \\
\hline
\end{tabularx}
\end{center}
\end{table}

\vspace{-1.35em}

\noindent
Support level: \textit{supported}

\vspace{0.3em}

\noindent
Examples: Buch{\zeroLinker}deckel, Buch\textbf{=er}regal, Kind\textbf{=er}jacke

\noindent
Counterexamples: Kind\textbf{es}alter, Mann\textbf{s}volk, Rind{\zeroLinker}fleisch

\vspace{0.3em}

\noindent
Sources: \citet[229]{eisenberg_2013}, \citet[99]{ortner_mueller-bollhagen_1991}

\vspace{0.35em}
\hrule




% 8
\vspace{0.4em}

\begin{center}
\textbf{Constraint (8)}
\end{center}

\noindent
First constituents that are constituted by nouns that build the plural form with -(")er attach a zero linker in majority of compounds in which the second constituent forces a singular meaning of the given first constituent.
\vspace{0.5em}

\begin{tikzpicture}
\node[
    draw,
    rounded corners=6pt,
    fill=gray!30,
    inner sep=2pt
] (11) {
    \colorbox{gray!30}{P2L}
};
\node[
    draw,
    rounded corners=6pt,
    fill=gray!30,
    inner sep=2pt,
    right=0.5em of 11
] (12) {
    \colorbox{gray!30}{MS}
};
\node[
    draw,
    rounded corners=6pt,
    fill=gray!30,
    inner sep=2pt,
    right=0.5em of 12
] {
    \colorbox{gray!30}{first constituent}
};
\end{tikzpicture}

\vspace{-0.25em}

\begin{table}[!ht]
\begin{center}
\setlength{\extrarowheight}{1pt}
\begin{tabularx}{\columnwidth}{|X|X|X|}
\hline
 & item & type \\
\hline
 coverage & $NA$ & $NA$ \\
\hline
 regularity & $NA$ & $NA$ \\
\hline
\end{tabularx}
\end{center}
\end{table}

\vspace{-1.35em}

\noindent
Support level: \textit{NA}

\vspace{0.3em}

\noindent
Examples: Buch{\zeroLinker}deckel, Grab{\zeroLinker}stein, Blatt{\zeroLinker}fläche

\noindent
Counterexamples: Mann\textbf{=er}herz, Kind\textbf{=er}bild, Buch{\zeroLinker}handel

\vspace{0.3em}

\noindent
Sources: \citet[229]{eisenberg_2013}, \citet[99, 108]{ortner_mueller-bollhagen_1991}

\vspace{0.35em}
\hrule




% 9
\vspace{0.4em}

\begin{center}
\textbf{Constraint (9)}
\end{center}

\noindent
First constituents that are constituted by nouns that build the plural form with -(")er tend to attach a zero linker in compounds in which the second constituent forces a mass meaning of the given first constituent.
\vspace{0.5em}

\begin{tikzpicture}
\node[
    draw,
    rounded corners=6pt,
    fill=gray!30,
    inner sep=2pt
] (11) {
    \colorbox{gray!30}{P2L}
};
\node[
    draw,
    rounded corners=6pt,
    fill=gray!30,
    inner sep=2pt,
    right=0.5em of 11
] (12) {
    \colorbox{gray!30}{MS}
};
\node[
    draw,
    rounded corners=6pt,
    fill=gray!30,
    inner sep=2pt,
    right=0.5em of 12
] {
    \colorbox{gray!30}{first constituent}
};
\end{tikzpicture}

\vspace{-0.25em}

\begin{table}[!ht]
\begin{center}
\setlength{\extrarowheight}{1pt}
\begin{tabularx}{\columnwidth}{|X|X|X|}
\hline
 & item & type \\
\hline
 coverage & $NA$ & $NA$ \\
\hline
 regularity & $NA$ & $NA$ \\
\hline
\end{tabularx}
\end{center}
\end{table}

\vspace{-1.35em}

\noindent
Support level: \textit{NA}

\vspace{0.3em}

\noindent
Examples: Glas{\zeroLinker}auge, Kraut{\zeroLinker}salat, Horn{\zeroLinker}haut

\noindent
Counterexamples: Kraut{\zeroLinker}garten

\vspace{0.3em}

\noindent
Sources: \citet[81]{nuebling_szczepaniak_2013}

\vspace{0.35em}
\hrule




% 10
\vspace{0.4em}

\begin{center}
\textbf{Constraint (10)}
\end{center}

\noindent
First constituents that are constituted by nouns that build the plural form with -(")er attach -(")er- more likely in compounds in which the second constituent forces a plural meaning of the given first constituent.
\vspace{0.5em}

\begin{tikzpicture}
\node[
    draw,
    rounded corners=6pt,
    fill=gray!30,
    inner sep=2pt
] (11) {
    \colorbox{gray!30}{P2L}
};
\node[
    draw,
    rounded corners=6pt,
    fill=gray!30,
    inner sep=2pt,
    right=0.5em of 11
] (12) {
    \colorbox{gray!30}{MS}
};
\node[
    draw,
    rounded corners=6pt,
    fill=gray!30,
    inner sep=2pt,
    right=0.5em of 12
] {
    \colorbox{gray!30}{first constituent}
};
\end{tikzpicture}

\vspace{-0.25em}

\begin{table}[!ht]
\begin{center}
\setlength{\extrarowheight}{1pt}
\begin{tabularx}{\columnwidth}{|X|X|X|}
\hline
 & item & type \\
\hline
 coverage & $NA$ & $NA$ \\
\hline
 regularity & $NA$ & $NA$ \\
\hline
\end{tabularx}
\end{center}
\end{table}

\vspace{-1.35em}

\noindent
Support level: \textit{NA}

\vspace{0.3em}

\noindent
Examples: Buch\textbf{=er}regal, Grab\textbf{=er}feld, Kraut\textbf{=er}frau

\noindent
Counterexamples: Buch{\zeroLinker}handel, Bild{\zeroLinker}band, Kind\textbf{=er}jacke, Mann\textbf{=er}herz

\vspace{0.3em}

\noindent
Sources: \citet[229]{eisenberg_2013}, \citet[80, 81]{nuebling_szczepaniak_2013}, \citet[99]{ortner_mueller-bollhagen_1991}, \citet[29]{schaefer_pankratz_2018}

\vspace{0.35em}
\hrule




% 11
\vspace{0.4em}

\begin{center}
\textbf{Constraint (11)}
\end{center}

\noindent
First constituents that are constituted by person and animal designations that build the plural form with -(")er may attach -(")er- regardless of the singular/plural interpretation of the given first constituent within the compound.
\vspace{0.5em}

\begin{tikzpicture}
\node[
    draw,
    rounded corners=6pt,
    fill=gray!30,
    inner sep=2pt
] (11) {
    \colorbox{gray!30}{P2L}
};
\node[
    draw,
    rounded corners=6pt,
    fill=gray!30,
    inner sep=2pt,
    right=0.5em of 11
] (12) {
    \colorbox{gray!30}{MS}
};
\node[
    draw,
    rounded corners=6pt,
    fill=gray!30,
    inner sep=2pt,
    right=0.5em of 12
] {
    \colorbox{gray!30}{first constituent}
};
\end{tikzpicture}

\vspace{-0.25em}

\begin{table}[!ht]
\begin{center}
\setlength{\extrarowheight}{1pt}
\begin{tabularx}{\columnwidth}{|X|X|X|}
\hline
 & item & type \\
\hline
 coverage & $NA$ & $NA$ \\
\hline
 regularity & $NA$ & $NA$ \\
\hline
\end{tabularx}
\end{center}
\end{table}

\vspace{-1.35em}

\noindent
Support level: \textit{NA}

\vspace{0.3em}

\noindent
Examples: Kind\textbf{=er}jacke, Huhn\textbf{=er}ei, Rind\textbf{=er}brust

\noindent
Counterexamples: 

\vspace{0.3em}

\noindent
Sources: \citet[229]{eisenberg_2013}, \citet[98]{ortner_mueller-bollhagen_1991}

\vspace{0.35em}
\hrule




% 12
\vspace{0.4em}

\begin{center}
\textbf{Constraint (12)}
\end{center}

\noindent
First constituents that are constituted by nouns that build the plural form with -e and umlaut mostly have a zero linker.
\vspace{0.5em}

\begin{tikzpicture}
\node[
    draw,
    rounded corners=6pt,
    fill=gray!30,
    inner sep=2pt
] (11) {
    \colorbox{gray!30}{P2L}
};
\node[
    draw,
    rounded corners=6pt,
    fill=gray!30,
    inner sep=2pt,
    right=0.5em of 11
] (12) {
    \colorbox{gray!30}{M}
};
\node[
    draw,
    rounded corners=6pt,
    fill=gray!30,
    inner sep=2pt,
    right=0.5em of 12
] {
    \colorbox{gray!30}{first constituent}
};
\end{tikzpicture}

\vspace{-0.25em}

\begin{table}[!ht]
\begin{center}
\setlength{\extrarowheight}{1pt}
\begin{tabularx}{\columnwidth}{|X|X|X|}
\hline
 & item & type \\
\hline
 coverage & $0.08$ & $0.111$ \\
\hline
 regularity & $0.722$ & $0.756$ \\
\hline
\end{tabularx}
\end{center}
\end{table}

\vspace{-1.35em}

\noindent
Support level: \textit{supported}

\vspace{0.3em}

\noindent
Examples: Hand{\zeroLinker}fläche, Frucht{\zeroLinker}joghurt, Arzt{\zeroLinker}praxis

\noindent
Counterexamples: Hand\textbf{=e}druck, Arzt\textbf{=e}streik, Frucht\textbf{=e}tee

\vspace{0.3em}

\noindent
Sources: \citet[83, 107]{ortner_mueller-bollhagen_1991}

\vspace{0.35em}
\hrule




% 13
\vspace{0.4em}

\begin{center}
\textbf{Constraint (13)}
\end{center}

\noindent
First constituents that are constituted by nouns that build the plural form with -e and umlaut attach -"e- irregularly in compounds in which the second constituent forces a plural meaning of the given first constituent.
\vspace{0.5em}

\begin{tikzpicture}
\node[
    draw,
    rounded corners=6pt,
    fill=gray!30,
    inner sep=2pt
] (11) {
    \colorbox{gray!30}{P2L}
};
\node[
    draw,
    rounded corners=6pt,
    fill=gray!30,
    inner sep=2pt,
    right=0.5em of 11
] (12) {
    \colorbox{gray!30}{MS}
};
\node[
    draw,
    rounded corners=6pt,
    fill=gray!30,
    inner sep=2pt,
    right=0.5em of 12
] {
    \colorbox{gray!30}{first constituent}
};
\end{tikzpicture}

\vspace{-0.25em}

\begin{table}[!ht]
\begin{center}
\setlength{\extrarowheight}{1pt}
\begin{tabularx}{\columnwidth}{|X|X|X|}
\hline
 & item & type \\
\hline
 coverage & $NA$ & $NA$ \\
\hline
 regularity & $NA$ & $NA$ \\
\hline
\end{tabularx}
\end{center}
\end{table}

\vspace{-1.35em}

\noindent
Support level: \textit{NA}

\vspace{0.3em}

\noindent
Examples: Hand\textbf{=e}druck, Arzt\textbf{=e}streik, Gast\textbf{=e}buch

\noindent
Counterexamples: Baum{\zeroLinker}gruppe, Zug{\zeroLinker}verkehr

\vspace{0.3em}

\noindent
Sources: \citet[81]{nuebling_szczepaniak_2013}, \citet[29]{schaefer_pankratz_2018}

\vspace{0.35em}
\hrule




% 14
\vspace{0.4em}

\begin{center}
\textbf{Constraint (14)}
\end{center}

\noindent
First constituents that are constituted by nouns that build the plural form with a zero ending and umlaut almost always have a zero linker.
\vspace{0.5em}

\begin{tikzpicture}
\node[
    draw,
    rounded corners=6pt,
    fill=gray!30,
    inner sep=2pt
] (11) {
    \colorbox{gray!30}{P2L}
};
\node[
    draw,
    rounded corners=6pt,
    fill=gray!30,
    inner sep=2pt,
    right=0.5em of 11
] (12) {
    \colorbox{gray!30}{M}
};
\node[
    draw,
    rounded corners=6pt,
    fill=gray!30,
    inner sep=2pt,
    right=0.5em of 12
] {
    \colorbox{gray!30}{first constituent}
};
\end{tikzpicture}

\vspace{-0.25em}

\begin{table}[!ht]
\begin{center}
\setlength{\extrarowheight}{1pt}
\begin{tabularx}{\columnwidth}{|X|X|X|}
\hline
 & item & type \\
\hline
 coverage & $0.007$ & $0.016$ \\
\hline
 regularity & $0.946$ & $0.967$ \\
\hline
\end{tabularx}
\end{center}
\end{table}

\vspace{-1.35em}

\noindent
Support level: \textit{supported}

\vspace{0.3em}

\noindent
Examples: Apfel{\zeroLinker}baum, Kloster{\zeroLinker}kirche, Mutter{\zeroLinker}sprache

\noindent
Counterexamples: Bruder\textbf{=}gemeinde, Mutter\textbf{=}zentrum, Vater\textbf{=}aufbruch

\vspace{0.3em}

\noindent
Sources: \citet[83, 106]{ortner_mueller-bollhagen_1991}

\vspace{0.35em}
\hrule




% 15
\vspace{0.4em}

\begin{center}
\textbf{Constraint (15)}
\end{center}

\noindent
First constituents that are constituted by nouns that build the plural form with a zero ending and umlaut attach a zero linker with umlaut irregularly in compounds in which the second constituent forces a plural meaning of the given first constituent.
\vspace{0.5em}

\begin{tikzpicture}
\node[
    draw,
    rounded corners=6pt,
    fill=gray!30,
    inner sep=2pt
] (11) {
    \colorbox{gray!30}{P2L}
};
\node[
    draw,
    rounded corners=6pt,
    fill=gray!30,
    inner sep=2pt,
    right=0.5em of 11
] (12) {
    \colorbox{gray!30}{MS}
};
\node[
    draw,
    rounded corners=6pt,
    fill=gray!30,
    inner sep=2pt,
    right=0.5em of 12
] {
    \colorbox{gray!30}{first constituent}
};
\end{tikzpicture}

\vspace{-0.25em}

\begin{table}[!ht]
\begin{center}
\setlength{\extrarowheight}{1pt}
\begin{tabularx}{\columnwidth}{|X|X|X|}
\hline
 & item & type \\
\hline
 coverage & $NA$ & $NA$ \\
\hline
 regularity & $NA$ & $NA$ \\
\hline
\end{tabularx}
\end{center}
\end{table}

\vspace{-1.35em}

\noindent
Support level: \textit{NA}

\vspace{0.3em}

\noindent
Examples: Bruder\textbf{=}gemeinde, Mutter\textbf{=}zentrum, Vater\textbf{=}aufbruch

\noindent
Counterexamples: Vogel{\zeroLinker}futter

\vspace{0.3em}

\noindent
Sources: \citet[29]{schaefer_pankratz_2018}

\vspace{0.35em}
\hrule




% 16
\vspace{0.4em}

\begin{center}
\textbf{Constraint (16)}
\end{center}

\noindent
First constituents that are constituted by mixed masculine and neuter nouns almost always attach -s-, a zero linker, or -(e)n-.
\vspace{0.5em}

\begin{tikzpicture}
\node[
    draw,
    rounded corners=6pt,
    fill=gray!30,
    inner sep=2pt
] (11) {
    \colorbox{gray!30}{P2L}
};
\node[
    draw,
    rounded corners=6pt,
    fill=gray!30,
    inner sep=2pt,
    right=0.5em of 11
] (12) {
    \colorbox{gray!30}{M}
};
\node[
    draw,
    rounded corners=6pt,
    fill=gray!30,
    inner sep=2pt,
    right=0.5em of 12
] {
    \colorbox{gray!30}{first constituent}
};
\end{tikzpicture}

\vspace{-0.25em}

\begin{table}[!ht]
\begin{center}
\setlength{\extrarowheight}{1pt}
\begin{tabularx}{\columnwidth}{|X|X|X|}
\hline
 & item & type \\
\hline
 coverage & $0.013$ & $0.013$ \\
\hline
 regularity & $0.997$ & $0.996$ \\
\hline
\end{tabularx}
\end{center}
\end{table}

\vspace{-1.35em}

\noindent
Support level: \textit{supported}

\vspace{0.3em}

\noindent
Examples: Staat\textbf{s}amt, Staat\textbf{en}bund, Bett{\zeroLinker}anzug, Auge\textbf{n}lied

\noindent
Counterexamples: Schmerz\textbf{ens}schrei

\vspace{0.3em}

\noindent
Sources: \citet[3]{krott_etal_2007}, \citet[80]{nuebling_szczepaniak_2013}, \citet[92]{ortner_mueller-bollhagen_1991}, \citet[229]{schaefer_2018}, \citet[9]{schaefer_pankratz_2018}

\vspace{0.35em}
\hrule




% 17
\vspace{0.4em}

\begin{center}
\textbf{Constraint (17)}
\end{center}

\noindent
First constituents that are constituted by mixed masculine and neuter nouns attach -(e)n- more likely in compounds in which the second constituent forces a plural meaning of the given first constituent.
\vspace{0.5em}

\begin{tikzpicture}
\node[
    draw,
    rounded corners=6pt,
    fill=gray!30,
    inner sep=2pt
] (11) {
    \colorbox{gray!30}{P2L}
};
\node[
    draw,
    rounded corners=6pt,
    fill=gray!30,
    inner sep=2pt,
    right=0.5em of 11
] (12) {
    \colorbox{gray!30}{MS}
};
\node[
    draw,
    rounded corners=6pt,
    fill=gray!30,
    inner sep=2pt,
    right=0.5em of 12
] {
    \colorbox{gray!30}{first constituent}
};
\end{tikzpicture}

\vspace{-0.25em}

\begin{table}[!ht]
\begin{center}
\setlength{\extrarowheight}{1pt}
\begin{tabularx}{\columnwidth}{|X|X|X|}
\hline
 & item & type \\
\hline
 coverage & $NA$ & $NA$ \\
\hline
 regularity & $NA$ & $NA$ \\
\hline
\end{tabularx}
\end{center}
\end{table}

\vspace{-1.35em}

\noindent
Support level: \textit{NA}

\vspace{0.3em}

\noindent
Examples: Staat\textbf{en}bund, Bett\textbf{en}zahl, Strahl\textbf{en}belastung

\noindent
Counterexamples: Motor\textbf{en}geräusch, Professor\textbf{en}titel, Auge\textbf{n}lied, Interesse\textbf{n}bereich

\vspace{0.3em}

\noindent
Sources: \citet[80]{nuebling_szczepaniak_2013}, \citet[92]{ortner_mueller-bollhagen_1991}

\vspace{0.35em}
\hrule




% 18
\vspace{0.4em}

\begin{center}
\textbf{Constraint (18)}
\end{center}

\noindent
First constituents that are constituted by derived nouns with suffixes -er, -ler, -ner, -el, -sel, -chen, -lein that build the plural form with a zero ending attach a zero linking element regularly.
\vspace{0.5em}

\begin{tikzpicture}
\node[
    draw,
    rounded corners=6pt,
    fill=gray!30,
    inner sep=2pt
] (11) {
    \colorbox{gray!30}{P2L}
};
\node[
    draw,
    rounded corners=6pt,
    fill=gray!30,
    inner sep=2pt,
    right=0.5em of 11
] (12) {
    \colorbox{gray!30}{D}
};
\node[
    draw,
    rounded corners=6pt,
    fill=gray!30,
    inner sep=2pt,
    right=0.5em of 12
] {
    \colorbox{gray!30}{first constituent}
};
\end{tikzpicture}

\vspace{-0.25em}

\begin{table}[!ht]
\begin{center}
\setlength{\extrarowheight}{1pt}
\begin{tabularx}{\columnwidth}{|X|X|X|}
\hline
 & item & type \\
\hline
 coverage & $0.071$ & $0.045$ \\
\hline
 regularity & $0.983$ & $0.968$ \\
\hline
\end{tabularx}
\end{center}
\end{table}

\vspace{-1.35em}

\noindent
Support level: \textit{supported}

\vspace{0.3em}

\noindent
Examples: Reiter{\zeroLinker}schwert, Gürtel{\zeroLinker}schnalle, Brötchen{\zeroLinker}geber

\noindent
Counterexamples: Alter\textbf{s}abstand, Reiter\textbf{s}mann

\vspace{0.3em}

\noindent
Sources: \citet[56, 82, 106]{ortner_mueller-bollhagen_1991}

\vspace{0.35em}
\hrule




% 19
\vspace{0.4em}

\begin{center}
\textbf{Constraint (19)}
\end{center}

\noindent
First constituents that are constituted by derived nouns with suffixes -bold, -nis, -rich, -at, -al that build the plural form with -e attach a zero linker regularly.
\vspace{0.5em}

\begin{tikzpicture}
\node[
    draw,
    rounded corners=6pt,
    fill=gray!30,
    inner sep=2pt
] (11) {
    \colorbox{gray!30}{P2L}
};
\node[
    draw,
    rounded corners=6pt,
    fill=gray!30,
    inner sep=2pt,
    right=0.5em of 11
] (12) {
    \colorbox{gray!30}{D}
};
\node[
    draw,
    rounded corners=6pt,
    fill=gray!30,
    inner sep=2pt,
    right=0.5em of 12
] {
    \colorbox{gray!30}{first constituent}
};
\end{tikzpicture}

\vspace{-0.25em}

\begin{table}[!ht]
\begin{center}
\setlength{\extrarowheight}{1pt}
\begin{tabularx}{\columnwidth}{|X|X|X|}
\hline
 & item & type \\
\hline
 coverage & $0.007$ & $0.005$ \\
\hline
 regularity & $0.747$ & $0.873$ \\
\hline
\end{tabularx}
\end{center}
\end{table}

\vspace{-1.35em}

\noindent
Support level: \textit{not supported}

\vspace{0.3em}

\noindent
Examples: Erlaubnis{\zeroLinker}schein, Format{\zeroLinker}vorlage, Personal{\zeroLinker}ausweis

\noindent
Counterexamples: Radikal\textbf{e}fänger, Internat\textbf{s}leitung, Magistrat\textbf{s}verfassung

\vspace{0.3em}

\noindent
Sources: \citet[83, 107]{ortner_mueller-bollhagen_1991}

\vspace{0.35em}
\hrule




% 20
\vspace{0.4em}

\begin{center}
\textbf{Constraint (20)}
\end{center}

\noindent
First constituents that are constituted by deverbal feminine nouns with a schwa suffix mostly attach a zero linker.
\vspace{0.5em}

\begin{tikzpicture}
\node[
    draw,
    rounded corners=6pt,
    fill=gray!30,
    inner sep=2pt
] (11) {
    \colorbox{gray!30}{P2L}
};
\node[
    draw,
    rounded corners=6pt,
    fill=gray!30,
    inner sep=2pt,
    right=0.5em of 11
] (12) {
    \colorbox{gray!30}{D}
};
\node[
    draw,
    rounded corners=6pt,
    fill=gray!30,
    inner sep=2pt,
    right=0.5em of 12
] {
    \colorbox{gray!30}{first constituent}
};
\end{tikzpicture}

\vspace{-0.25em}

\begin{table}[!ht]
\begin{center}
\setlength{\extrarowheight}{1pt}
\begin{tabularx}{\columnwidth}{|X|X|X|}
\hline
 & item & type \\
\hline
 coverage & $0.028$ & $0.028$ \\
\hline
 regularity & $0.579$ & $0.592$ \\
\hline
\end{tabularx}
\end{center}
\end{table}

\vspace{-1.35em}

\noindent
Support level: \textit{supported}

\vspace{0.3em}

\noindent
Examples: Abgabe{\zeroLinker}soll, Sorge{\zeroLinker}pflicht, Lage{\zeroLinker}plan

\noindent
Counterexamples: Abgabe\textbf{n}ordnung, Anzeige\textbf{n}blatt, Lüge\textbf{n}detektor

\vspace{0.3em}

\noindent
Sources: \citet[93]{ortner_müller-bollhagen_1991}

\vspace{0.35em}
\hrule




% 21
\vspace{0.4em}

\begin{center}
\textbf{Constraint (21)}
\end{center}

\noindent
First constituents that are constituted by deverbal feminine nouns with a schwa suffix tend to attach -(e)n- in compounds in which the second constituent forces a plural reading of the given first constituent.
\vspace{0.5em}

\begin{tikzpicture}
\node[
    draw,
    rounded corners=6pt,
    fill=gray!30,
    inner sep=2pt
] (11) {
    \colorbox{gray!30}{P2L}
};
\node[
    draw,
    rounded corners=6pt,
    fill=gray!30,
    inner sep=2pt,
    right=0.5em of 11
] (12) {
    \colorbox{gray!30}{DS}
};
\node[
    draw,
    rounded corners=6pt,
    fill=gray!30,
    inner sep=2pt,
    right=0.5em of 12
] {
    \colorbox{gray!30}{first constituent}
};
\end{tikzpicture}

\vspace{-0.25em}

\begin{table}[!ht]
\begin{center}
\setlength{\extrarowheight}{1pt}
\begin{tabularx}{\columnwidth}{|X|X|X|}
\hline
 & item & type \\
\hline
 coverage & $NA$ & $NA$ \\
\hline
 regularity & $NA$ & $NA$ \\
\hline
\end{tabularx}
\end{center}
\end{table}

\vspace{-1.35em}

\noindent
Support level: \textit{NA}

\vspace{0.3em}

\noindent
Examples: Abgabe\textbf{n}ordnung, Anzeige\textbf{n}blatt, Lüge\textbf{n}detektor

\noindent
Counterexamples: Einnahme{\zeroLinker}buch, Anzeige{\zeroLinker}tafel

\vspace{0.3em}

\noindent
Sources: \citet[93]{ortner_müller-bollhagen_1991}

\vspace{0.35em}
\hrule




% 22
\vspace{0.4em}

\begin{center}
\textbf{Constraint (22)}
\end{center}

\noindent
First constituents that are constituted by deadjectival feminine nouns with a schwa suffix almost always attach -(e)n- or a zero linker. Each of the two linkers is preferred in about the same number of cases.
\vspace{0.5em}

\begin{tikzpicture}
\node[
    draw,
    rounded corners=6pt,
    fill=gray!30,
    inner sep=2pt
] (11) {
    \colorbox{gray!30}{P2L}
};
\node[
    draw,
    rounded corners=6pt,
    fill=gray!30,
    inner sep=2pt,
    right=0.5em of 11
] (12) {
    \colorbox{gray!30}{D}
};
\node[
    draw,
    rounded corners=6pt,
    fill=gray!30,
    inner sep=2pt,
    right=0.5em of 12
] {
    \colorbox{gray!30}{first constituent}
};
\end{tikzpicture}

\vspace{-0.25em}

\begin{table}[!ht]
\begin{center}
\setlength{\extrarowheight}{1pt}
\begin{tabularx}{\columnwidth}{|X|X|X|}
\hline
 & item & type \\
\hline
 coverage & $0.009$ & $0.01$ \\
\hline
 regularity & $0.97$ & $0.852$ \\
\hline
\end{tabularx}
\end{center}
\end{table}

\vspace{-1.35em}

\noindent
Support level: \textit{supported}

\vspace{0.3em}

\noindent
Examples: Fläche{\zeroLinker}maß, Hitze{\zeroLinker}welle, Tiefe\textbf{n}meter, Größe\textbf{n}klasse

\noindent
Counterexamples: Liebe\textbf{s}brief

\vspace{0.3em}

\noindent
Sources: \citet[573]{fuhrhop_kürschner_2015}, \citet[112]{kürschner_2005}, \citet[93]{ortner_müller-bollhagen_1991}

\vspace{0.35em}
\hrule




% 23
\vspace{0.4em}

\begin{center}
\textbf{Constraint (23)}
\end{center}

\noindent
First constituents that are constituted by derived nouns with suffixes -(ig)keit, -heit, -schaft, -ung, -sal, -ing, -ling, -tum, -um, -ion, also -ität and its allomorphs attach -s- regularly.
\vspace{0.5em}

\begin{tikzpicture}
\node[
    draw,
    rounded corners=6pt,
    fill=gray!30,
    inner sep=2pt
] (11) {
    \colorbox{gray!30}{P2L}
};
\node[
    draw,
    rounded corners=6pt,
    fill=gray!30,
    inner sep=2pt,
    right=0.5em of 11
] (12) {
    \colorbox{gray!30}{D}
};
\node[
    draw,
    rounded corners=6pt,
    fill=gray!30,
    inner sep=2pt,
    right=0.5em of 12
] {
    \colorbox{gray!30}{first constituent}
};
\end{tikzpicture}

\vspace{-0.25em}

\begin{table}[!ht]
\begin{center}
\setlength{\extrarowheight}{1pt}
\begin{tabularx}{\columnwidth}{|X|X|X|}
\hline
 & item & type \\
\hline
 coverage & $0.215$ & $0.158$ \\
\hline
 regularity & $0.98$ & $0.985$ \\
\hline
\end{tabularx}
\end{center}
\end{table}

\vspace{-1.35em}

\noindent
Support level: \textit{supported}

\vspace{0.3em}

\noindent
Examples: Gesundheit\textbf{s}amt, Gesellschaft\textbf{s}politik, Schicksal\textbf{s}drama, Rarität\textbf{s}wert, Eigentum\textbf{s}recht

\noindent
Counterexamples: Minderheit\textbf{en}recht, Kuriosität\textbf{en}museum, Nationalität\textbf{en}kampf

\vspace{0.3em}

\noindent
Sources: \citet[230]{eisenberg_2013}, \citet[572]{fuhrhop_kuerschner_2015}, \citet[61]{koliopoulou_2014}, \citet[3]{krott_etal_2007}, \citet[107, 115]{kuerschner_2005}, \citet[20, 21]{kuerschner_2010}, \citet[10, 20, 23]{nuebling_szczepaniak_2008}, \citet[77]{nuebling_szczepaniak_2013}, \citet[73, 83, 88, 89, 94]{ortner_mueller-bollhagen_1991}, \citet[229]{schaefer_2018}, \citet[8]{schaefer_pankratz_2018}, \citet[19]{schluecker_2023}, \citet[448]{wegener_2003}

\vspace{0.35em}
\hrule




% 24
\vspace{0.4em}

\begin{center}
\textbf{Constraint (24)}
\end{center}

\noindent
First constituents that are constituted by derived nouns with suffix -ität and its allomorphs prefer -en- in compounds in which the second constituent forces a plural meaning of the given first constituent.
\vspace{0.5em}

\begin{tikzpicture}
\node[
    draw,
    rounded corners=6pt,
    fill=gray!30,
    inner sep=2pt
] (11) {
    \colorbox{gray!30}{P2L}
};
\node[
    draw,
    rounded corners=6pt,
    fill=gray!30,
    inner sep=2pt,
    right=0.5em of 11
] (12) {
    \colorbox{gray!30}{DS}
};
\node[
    draw,
    rounded corners=6pt,
    fill=gray!30,
    inner sep=2pt,
    right=0.5em of 12
] {
    \colorbox{gray!30}{first constituent}
};
\end{tikzpicture}

\vspace{-0.25em}

\begin{table}[!ht]
\begin{center}
\setlength{\extrarowheight}{1pt}
\begin{tabularx}{\columnwidth}{|X|X|X|}
\hline
 & item & type \\
\hline
 coverage & $NA$ & $NA$ \\
\hline
 regularity & $NA$ & $NA$ \\
\hline
\end{tabularx}
\end{center}
\end{table}

\vspace{-1.35em}

\noindent
Support level: \textit{NA}

\vspace{0.3em}

\noindent
Examples: Rarität\textbf{en}sammlung, Kuriosität\textbf{en}händler, Aktualitäten\textbf{en}kino

\noindent
Counterexamples: 

\vspace{0.3em}

\noindent
Sources: \citet[94]{ortner_mueller-bollhagen_1991}

\vspace{0.35em}
\hrule




% 25
\vspace{0.4em}

\begin{center}
\textbf{Constraint (25)}
\end{center}

\noindent
First constituents that are constituted by deverbal nouns that end in suffix -en attach -s- regularly.
\vspace{0.5em}

\begin{tikzpicture}
\node[
    draw,
    rounded corners=6pt,
    fill=gray!30,
    inner sep=2pt
] (11) {
    \colorbox{gray!30}{P2L}
};
\node[
    draw,
    rounded corners=6pt,
    fill=gray!30,
    inner sep=2pt,
    right=0.5em of 11
] (12) {
    \colorbox{gray!30}{D}
};
\node[
    draw,
    rounded corners=6pt,
    fill=gray!30,
    inner sep=2pt,
    right=0.5em of 12
] {
    \colorbox{gray!30}{first constituent}
};
\end{tikzpicture}

\vspace{-0.25em}

\begin{table}[!ht]
\begin{center}
\setlength{\extrarowheight}{1pt}
\begin{tabularx}{\columnwidth}{|X|X|X|}
\hline
 & item & type \\
\hline
 coverage & $0.005$ & $0.009$ \\
\hline
 regularity & $0.773$ & $0.926$ \\
\hline
\end{tabularx}
\end{center}
\end{table}

\vspace{-1.35em}

\noindent
Support level: \textit{not supported}

\vspace{0.3em}

\noindent
Examples: Leben\textbf{s}mittel, Essen\textbf{s}gewohnheit, Unternehmen\textbf{s}bereich

\noindent
Counterexamples: Essen{\zeroLinker}portion, Leben{\zeroLinker}geschäft, Husten{\zeroLinker}tropfen, Beben{\zeroLinker}gebiet

\vspace{0.3em}

\noindent
Sources: \citet[230]{eisenberg_2013}, \citet[107]{kuerschner_2005}, \citet[78]{nuebling_szczepaniak_2013}, \citet[89]{ortner_mueller-bollhagen_1991}

\vspace{0.35em}
\hrule




% 26
\vspace{0.4em}

\begin{center}
\textbf{Constraint (26)}
\end{center}

\noindent
First constituents that are constituted by derived feminine nouns with suffix -in always attach -en-. (The suffix -in adjusts orthographically in this case and becomes -inn.)
\vspace{0.5em}

\begin{tikzpicture}
\node[
    draw,
    rounded corners=6pt,
    fill=gray!30,
    inner sep=2pt
] (11) {
    \colorbox{gray!30}{P2L}
};
\node[
    draw,
    rounded corners=6pt,
    fill=gray!30,
    inner sep=2pt,
    right=0.5em of 11
] (12) {
    \colorbox{gray!30}{D}
};
\node[
    draw,
    rounded corners=6pt,
    fill=gray!30,
    inner sep=2pt,
    right=0.5em of 12
] {
    \colorbox{gray!30}{first constituent}
};
\end{tikzpicture}

\vspace{-0.25em}

\begin{table}[!ht]
\begin{center}
\setlength{\extrarowheight}{1pt}
\begin{tabularx}{\columnwidth}{|X|X|X|}
\hline
 & item & type \\
\hline
 coverage & $NA$ & $NA$ \\
\hline
 regularity & $NA$ & $NA$ \\
\hline
\end{tabularx}
\end{center}
\end{table}

\vspace{-1.35em}

\noindent
Support level: \textit{NA}

\vspace{0.3em}

\noindent
Examples: Lehrerinn\textbf{en}mentalität, Freundinn\textbf{en}gruppe, Schülerinn\textbf{en}rat

\noindent
Counterexamples: Königin{\zeroLinker}mutter

\vspace{0.3em}

\noindent
Sources: \citet[94]{ortner_mueller-bollhagen_1991}

\vspace{0.35em}
\hrule




% 27
\vspace{0.4em}

\begin{center}
\textbf{Constraint (27)}
\end{center}

\noindent
First constituents that are constituted by prefixed deverbal nouns exhibit a tendency to attach -s-.
\vspace{0.5em}

\begin{tikzpicture}
\node[
    draw,
    rounded corners=6pt,
    fill=gray!30,
    inner sep=2pt
] (11) {
    \colorbox{gray!30}{P2L}
};
\node[
    draw,
    rounded corners=6pt,
    fill=gray!30,
    inner sep=2pt,
    right=0.5em of 11
] (12) {
    \colorbox{gray!30}{D}
};
\node[
    draw,
    rounded corners=6pt,
    fill=gray!30,
    inner sep=2pt,
    right=0.5em of 12
] {
    \colorbox{gray!30}{first constituent}
};
\end{tikzpicture}

\vspace{-0.25em}

\begin{table}[!ht]
\begin{center}
\setlength{\extrarowheight}{1pt}
\begin{tabularx}{\columnwidth}{|X|X|X|}
\hline
 & item & type \\
\hline
 coverage & $0.159$ & $0.122$ \\
\hline
 regularity & $0.641$ & $0.659$ \\
\hline
\end{tabularx}
\end{center}
\end{table}

\vspace{-1.35em}

\noindent
Support level: \textit{supported}

\vspace{0.3em}

\noindent
Examples: Anspruch\textbf{s}haltung, Eintrag\textbf{s}frist, Bedarf\textbf{s}fall, Verfall\textbf{s}datum

\noindent
Counterexamples: Anruf{\zeroLinker}beantworter, Überfall{\zeroLinker}kommando, Bestand{\zeroLinker}teil

\vspace{0.3em}

\noindent
Sources: \citet[230]{eisenberg_2013}, \citet[61]{koliopoulou_2014}, \citet[107]{kuerschner_2005}, \citet[21]{kuerschner_2010}, \citet[190]{kopf_2017}, \citet[18, 19]{nuebling_szczepaniak_2008}, \citet[78]{nuebling_szczepaniak_2013}, \citet[89]{ortner_mueller-bollhagen_1991}

\vspace{0.35em}
\hrule




% 28
\vspace{0.4em}

\begin{center}
\textbf{Constraint (28)}
\end{center}

\noindent
First constituents that are constituted by nouns that end in a sibilant or in a consonant cluster including [s] mostly adopt a zero linker.
\vspace{0.5em}

\begin{tikzpicture}
\node[
    draw,
    rounded corners=6pt,
    fill=gray!30,
    inner sep=2pt
] (11) {
    \colorbox{gray!30}{P2L}
};
\node[
    draw,
    rounded corners=6pt,
    fill=gray!30,
    inner sep=2pt,
    right=0.5em of 11
] (12) {
    \colorbox{gray!30}{P}
};
\node[
    draw,
    rounded corners=6pt,
    fill=gray!30,
    inner sep=2pt,
    right=0.5em of 12
] {
    \colorbox{gray!30}{first constituent}
};
\end{tikzpicture}

\vspace{-0.25em}

\begin{table}[!ht]
\begin{center}
\setlength{\extrarowheight}{1pt}
\begin{tabularx}{\columnwidth}{|X|X|X|}
\hline
 & item & type \\
\hline
 coverage & $0.083$ & $0.1$ \\
\hline
 regularity & $0.97$ & $0.953$ \\
\hline
\end{tabularx}
\end{center}
\end{table}

\vspace{-1.35em}

\noindent
Support level: \textit{supported}

\vspace{0.3em}

\noindent
Examples: Gefäß{\zeroLinker}system, Fisch{\zeroLinker}öl, Notiz{\zeroLinker}buch, Herbst{\zeroLinker}anfang

\noindent
Counterexamples: Gast\textbf{=e}buch, Geist\textbf{=er}haus

\vspace{0.3em}

\noindent
Sources: \citet[19]{kuerschner_2010}, \citet[89, 109]{ortner_mueller-bollhagen_1991}, \citet[181]{wegener_2005}

\vspace{0.35em}
\hrule




% 29
\vspace{0.4em}

\begin{center}
\textbf{Constraint (29)}
\end{center}

\noindent
First constituents that are constituted by nouns that end in a full vowel always adopt a zero linker.
\vspace{0.5em}

\begin{tikzpicture}
\node[
    draw,
    rounded corners=6pt,
    fill=gray!30,
    inner sep=2pt
] (11) {
    \colorbox{gray!30}{P2L}
};
\node[
    draw,
    rounded corners=6pt,
    fill=gray!30,
    inner sep=2pt,
    right=0.5em of 11
] (12) {
    \colorbox{gray!30}{P}
};
\node[
    draw,
    rounded corners=6pt,
    fill=gray!30,
    inner sep=2pt,
    right=0.5em of 12
] {
    \colorbox{gray!30}{first constituent}
};
\end{tikzpicture}

\vspace{-0.25em}

\begin{table}[!ht]
\begin{center}
\setlength{\extrarowheight}{1pt}
\begin{tabularx}{\columnwidth}{|X|X|X|}
\hline
 & item & type \\
\hline
 coverage & $0.041$ & $0.036$ \\
\hline
 regularity & $0.968$ & $0.976$ \\
\hline
\end{tabularx}
\end{center}
\end{table}

\vspace{-1.35em}

\noindent
Support level: \textit{partially supported}

\vspace{0.3em}

\noindent
Examples: Auto{\zeroLinker}bahn, Uhu{\zeroLinker}paar, Kuh{\zeroLinker}milch, Gummi{\zeroLinker}bär

\noindent
Counterexamples: Uhu\textbf{s}nest, Idee\textbf{n}liste, Kalorie\textbf{n}wert, Schuh\textbf{e}kauf

\vspace{0.3em}

\noindent
Sources: \citet[110]{kuerschner_2005}, \citet[9]{schaefer_pankratz_2018}

\vspace{0.35em}
\hrule




% 30
\vspace{0.4em}

\begin{center}
\textbf{Constraint (30)}
\end{center}

\noindent
First constituents that are constituted by feminine nouns that end with stressed -ei, -ie, -ur, also stressed or unstressed -ik usually attach a zero linker.
\vspace{0.5em}

\begin{tikzpicture}
\node[
    draw,
    rounded corners=6pt,
    fill=gray!30,
    inner sep=2pt
] (11) {
    \colorbox{gray!30}{P2L}
};
\node[
    draw,
    rounded corners=6pt,
    fill=gray!30,
    inner sep=2pt,
    right=0.5em of 11
] (12) {
    \colorbox{gray!30}{MP}
};
\node[
    draw,
    rounded corners=6pt,
    fill=gray!30,
    inner sep=2pt,
    right=0.5em of 12
] {
    \colorbox{gray!30}{first constituent}
};
\end{tikzpicture}

\vspace{-0.25em}

\begin{table}[!ht]
\begin{center}
\setlength{\extrarowheight}{1pt}
\begin{tabularx}{\columnwidth}{|X|X|X|}
\hline
 & item & type \\
\hline
 coverage & $0.03$ & $0.033$ \\
\hline
 regularity & $0.95$ & $0.969$ \\
\hline
\end{tabularx}
\end{center}
\end{table}

\vspace{-1.35em}

\noindent
Support level: \textit{supported}

\vspace{0.3em}

\noindent
Examples: Kosmetik{\zeroLinker}industrie, Akustik{\zeroLinker}paneele, Architektur{\zeroLinker}büro, Metzgerei{\zeroLinker}produkt

\noindent
Counterexamples: Kultur\textbf{en}folge, Melodie\textbf{n}reigen, Partei\textbf{en}staat

\vspace{0.3em}

\noindent
Sources: \citet[107]{ortner_mueller-bollhagen_1991}

\vspace{0.35em}
\hrule




% 31
\vspace{0.4em}

\begin{center}
\textbf{Constraint (31)}
\end{center}

\noindent
First constituents that are constituted by polysyllabic feminine nouns that end with [t] often attach -s- if the [t] is not part of a suffix -(ig)keit, -heit, -schaft, or -ität or its allomorphs.
\vspace{0.5em}

\begin{tikzpicture}
\node[
    draw,
    rounded corners=6pt,
    fill=gray!30,
    inner sep=2pt
] (11) {
    \colorbox{gray!30}{P2L}
};
\node[
    draw,
    rounded corners=6pt,
    fill=gray!30,
    inner sep=2pt,
    right=0.5em of 11
] (12) {
    \colorbox{gray!30}{MDP}
};
\node[
    draw,
    rounded corners=6pt,
    fill=gray!30,
    inner sep=2pt,
    right=0.5em of 12
] {
    \colorbox{gray!30}{first constituent}
};
\end{tikzpicture}

\vspace{-0.25em}

\begin{table}[!ht]
\begin{center}
\setlength{\extrarowheight}{1pt}
\begin{tabularx}{\columnwidth}{|X|X|X|}
\hline
 & item & type \\
\hline
 coverage & $0.013$ & $0.016$ \\
\hline
 regularity & $0.505$ & $0.463$ \\
\hline
\end{tabularx}
\end{center}
\end{table}

\vspace{-1.35em}

\noindent
Support level: \textit{partially supported}

\vspace{0.3em}

\noindent
Examples: Arbeit\textbf{s}tag, Heirat\textbf{s}antrag, Zukunft\textbf{s}angst

\noindent
Counterexamples: Arbeit{\zeroLinker}geber, Jugend{\zeroLinker}alter, Umwelt{\zeroLinker}schutz

\vspace{0.3em}

\noindent
Sources: \citet[77]{nuebling_szczepaniak_2013}, \citet[74, 75]{ortner_mueller-bollhagen_1991}

\vspace{0.35em}
\hrule




% 32
\vspace{0.4em}

\begin{center}
\textbf{Constraint (32)}
\end{center}

\noindent
First constituents that are constituted by nouns that end in schwa mostly adopt -(e)n-.
\vspace{0.5em}

\begin{tikzpicture}
\node[
    draw,
    rounded corners=6pt,
    fill=gray!30,
    inner sep=2pt
] (11) {
    \colorbox{gray!30}{P2L}
};
\node[
    draw,
    rounded corners=6pt,
    fill=gray!30,
    inner sep=2pt,
    right=0.5em of 11
] (12) {
    \colorbox{gray!30}{P}
};
\node[
    draw,
    rounded corners=6pt,
    fill=gray!30,
    inner sep=2pt,
    right=0.5em of 12
] {
    \colorbox{gray!30}{first constituent}
};
\end{tikzpicture}

\vspace{-0.25em}

\begin{table}[!ht]
\begin{center}
\setlength{\extrarowheight}{1pt}
\begin{tabularx}{\columnwidth}{|X|X|X|}
\hline
 & item & type \\
\hline
 coverage & $0.172$ & $0.149$ \\
\hline
 regularity & $0.752$ & $0.72$ \\
\hline
\end{tabularx}
\end{center}
\end{table}

\vspace{-1.35em}

\noindent
Support level: \textit{supported}

\vspace{0.3em}

\noindent
Examples: Biene\textbf{n}zucht, Suppe\textbf{n}schüssel, Auge\textbf{n}lied

\noindent
Counterexamples: Fläche{\zeroLinker}maß, Sorge{\zeroLinker}pflicht, Gebäude{\zeroLinker}komplex, Gedanke\textbf{ns}turm

\vspace{0.3em}

\noindent
Sources: \citet[18]{kuerschner_2010}, \citet[83, 92]{ortner_mueller-bollhagen_1991}

\vspace{0.35em}
\hrule




% 33
\vspace{0.4em}

\begin{center}
\textbf{Constraint (33)}
\end{center}

\noindent
First constituents that are constituted by feminine nouns that end in schwa adopt -(e)n- regularly.
\vspace{0.5em}

\begin{tikzpicture}
\node[
    draw,
    rounded corners=6pt,
    fill=gray!30,
    inner sep=2pt
] (11) {
    \colorbox{gray!30}{P2L}
};
\node[
    draw,
    rounded corners=6pt,
    fill=gray!30,
    inner sep=2pt,
    right=0.5em of 11
] (12) {
    \colorbox{gray!30}{MP}
};
\node[
    draw,
    rounded corners=6pt,
    fill=gray!30,
    inner sep=2pt,
    right=0.5em of 12
] {
    \colorbox{gray!30}{first constituent}
};
\end{tikzpicture}

\vspace{-0.25em}

\begin{table}[!ht]
\begin{center}
\setlength{\extrarowheight}{1pt}
\begin{tabularx}{\columnwidth}{|X|X|X|}
\hline
 & item & type \\
\hline
 coverage & $0.156$ & $0.133$ \\
\hline
 regularity & $0.765$ & $0.744$ \\
\hline
\end{tabularx}
\end{center}
\end{table}

\vspace{-1.35em}

\noindent
Support level: \textit{not supported}

\vspace{0.3em}

\noindent
Examples: Biene\textbf{n}zucht, Suppe\textbf{n}schüssel, Tiefe\textbf{n}meter

\noindent
Counterexamples: Fläche{\zeroLinker}maß, Sorge{\zeroLinker}pflicht

\vspace{0.3em}

\noindent
Sources: \citet[229]{eisenberg_2013}, \citet[573]{fuhrhop_kuerschner_2015}, \citet[3]{krott_etal_2007}, \citet[151, 157]{libben_etal_2009}, \citet[83, 92]{ortner_mueller-bollhagen_1991}, \citet[9]{schaefer_pankratz_2018}

\vspace{0.35em}
\hrule




% 34
\vspace{0.4em}

\begin{center}
\textbf{Constraint (34)}
\end{center}

\noindent
First constituents that are constituted by consonant-final weak feminine nouns mostly adopt a zero linker or -s- in compounds in which the second constituent forces a singular meaning of the given first constituent.
\vspace{0.5em}

\begin{tikzpicture}
\node[
    draw,
    rounded corners=6pt,
    fill=gray!30,
    inner sep=2pt
] (11) {
    \colorbox{gray!30}{P2L}
};
\node[
    draw,
    rounded corners=6pt,
    fill=gray!30,
    inner sep=2pt,
    right=0.5em of 11
] (12) {
    \colorbox{gray!30}{MPS}
};
\node[
    draw,
    rounded corners=6pt,
    fill=gray!30,
    inner sep=2pt,
    right=0.5em of 12
] {
    \colorbox{gray!30}{first constituent}
};
\end{tikzpicture}

\vspace{-0.25em}

\begin{table}[!ht]
\begin{center}
\setlength{\extrarowheight}{1pt}
\begin{tabularx}{\columnwidth}{|X|X|X|}
\hline
 & item & type \\
\hline
 coverage & $NA$ & $NA$ \\
\hline
 regularity & $NA$ & $NA$ \\
\hline
\end{tabularx}
\end{center}
\end{table}

\vspace{-1.35em}

\noindent
Support level: \textit{NA}

\vspace{0.3em}

\noindent
Examples: Schrift{\zeroLinker}führer, Geburt\textbf{s}tag, Burg{\zeroLinker}anlage

\noindent
Counterexamples: Burg\textbf{en}blick, Frau\textbf{en}bild

\vspace{0.3em}

\noindent
Sources: \citet[229]{eisenberg_2013}, \citet[3]{krott_etal_2007}, \citet[80]{nuebling_szczepaniak_2013}, \citet[94, 110]{ortner_mueller-bollhagen_1991}

\vspace{0.35em}
\hrule




% 35
\vspace{0.4em}

\begin{center}
\textbf{Constraint (35)}
\end{center}

\noindent
First constituents that are constituted by consonant-final weak feminine nouns --- especially final-stressed incl. monosyllabic ones --- attach -en- more likely in compounds in which the second constituent forces a plural meaning of the given first constituent.
\vspace{0.5em}

\begin{tikzpicture}
\node[
    draw,
    rounded corners=6pt,
    fill=gray!30,
    inner sep=2pt
] (11) {
    \colorbox{gray!30}{P2L}
};
\node[
    draw,
    rounded corners=6pt,
    fill=gray!30,
    inner sep=2pt,
    right=0.5em of 11
] (12) {
    \colorbox{gray!30}{MPS}
};
\node[
    draw,
    rounded corners=6pt,
    fill=gray!30,
    inner sep=2pt,
    right=0.5em of 12
] {
    \colorbox{gray!30}{first constituent}
};
\end{tikzpicture}

\vspace{-0.25em}

\begin{table}[!ht]
\begin{center}
\setlength{\extrarowheight}{1pt}
\begin{tabularx}{\columnwidth}{|X|X|X|}
\hline
 & item & type \\
\hline
 coverage & $NA$ & $NA$ \\
\hline
 regularity & $NA$ & $NA$ \\
\hline
\end{tabularx}
\end{center}
\end{table}

\vspace{-1.35em}

\noindent
Support level: \textit{NA}

\vspace{0.3em}

\noindent
Examples: Schrift\textbf{en}verzeichnis, Geburt\textbf{en}kontrolle, Burg\textbf{en}land

\noindent
Counterexamples: Burg\textbf{en}blick

\vspace{0.3em}

\noindent
Sources: \citet[229]{eisenberg_2013}, \citet[3]{nuebling_szczepaniak_2008}, \citet[80]{nuebling_szczepaniak_2013}, \citet[110]{ortner_mueller-bollhagen_1991}, \citet[29]{schaefer_2018}

\vspace{0.35em}
\hrule




% 36
\vspace{0.4em}

\begin{center}
\textbf{Constraint (36)}
\end{center}

\noindent
Copulative compounds insert a zero linker regularly. By copulative compounds are understood compounds in which the two constituents do not exhibit a clear modifier-head relation.
\vspace{0.5em}

\begin{tikzpicture}
\node[
    draw,
    rounded corners=6pt,
    fill=gray!30,
    inner sep=2pt
] (11) {
    \colorbox{gray!30}{P2L}
};
\node[
    draw,
    rounded corners=6pt,
    fill=gray!30,
    inner sep=2pt,
    right=0.5em of 11
] (12) {
    \colorbox{gray!30}{S}
};
\node[
    draw,
    rounded corners=6pt,
    fill=gray!30,
    inner sep=2pt,
    right=0.5em of 12
] {
    \colorbox{gray!30}{first + second constituent}
};
\end{tikzpicture}

\vspace{-0.25em}

\begin{table}[!ht]
\begin{center}
\setlength{\extrarowheight}{1pt}
\begin{tabularx}{\columnwidth}{|X|X|X|}
\hline
 & item & type \\
\hline
 coverage & $NA$ & $NA$ \\
\hline
 regularity & $NA$ & $NA$ \\
\hline
\end{tabularx}
\end{center}
\end{table}

\vspace{-1.35em}

\noindent
Support level: \textit{NA}

\vspace{0.3em}

\noindent
Examples: Dichter{\zeroLinker}componist, Bett{\zeroLinker}sofa, Königin{\zeroLinker}mutter

\noindent
Counterexamples: 

\vspace{0.3em}

\noindent
Sources: \citet[63, 64]{koliopoulou_2014}, \citet[31, 32]{neef_borgwaldt_2012}, \citet[57, 91, 94, 109]{ortner_mueller-bollhagen_1991}

\vspace{0.35em}
\hrule




% 37
\vspace{0.4em}

\begin{center}
\textbf{Constraint (37)}
\end{center}

\noindent
Argumental compounds are sometimes marked by -s-. By argumental compounds are understood compounds in which the second constituent still contains a high degree of verbiness and the first constituent of which constitutes their argument. The second constituent is such compounds is usually an agent designation, an -ung formation etc.
\vspace{0.5em}

\begin{tikzpicture}
\node[
    draw,
    rounded corners=6pt,
    fill=gray!30,
    inner sep=2pt
] (11) {
    \colorbox{gray!30}{P2L}
};
\node[
    draw,
    rounded corners=6pt,
    fill=gray!30,
    inner sep=2pt,
    right=0.5em of 11
] (12) {
    \colorbox{gray!30}{S}
};
\node[
    draw,
    rounded corners=6pt,
    fill=gray!30,
    inner sep=2pt,
    right=0.5em of 12
] {
    \colorbox{gray!30}{first + second constituent}
};
\end{tikzpicture}

\vspace{-0.25em}

\begin{table}[!ht]
\begin{center}
\setlength{\extrarowheight}{1pt}
\begin{tabularx}{\columnwidth}{|X|X|X|}
\hline
 & item & type \\
\hline
 coverage & $NA$ & $NA$ \\
\hline
 regularity & $NA$ & $NA$ \\
\hline
\end{tabularx}
\end{center}
\end{table}

\vspace{-1.35em}

\noindent
Support level: \textit{NA}

\vspace{0.3em}

\noindent
Examples: Gewicht\textbf{s}heber, Krieg\textbf{s}führung, Unternehmen\textbf{s}beratung

\noindent
Counterexamples: Gewicht{\zeroLinker}heber, Krieg{\zeroLinker}führung

\vspace{0.3em}

\noindent
Sources: \citet[79]{nuebling_szczepaniak_2013}

\vspace{0.35em}
\hrule




% 38
\vspace{0.4em}

\begin{center}
\textbf{Constraint (38)}
\end{center}

\noindent
Compounds belonging to technical terminology (economics, law, medicine, etc.) are often missing an explicit linking element (even in cases where it would be obligatory in a non-technical context).
\vspace{0.5em}

\begin{tikzpicture}
\node[
    draw,
    rounded corners=6pt,
    fill=gray!30,
    inner sep=2pt
] (11) {
    \colorbox{gray!30}{P2L}
};
\node[
    draw,
    rounded corners=6pt,
    fill=gray!30,
    inner sep=2pt,
    right=0.5em of 11
] (12) {
    \colorbox{gray!30}{L}
};
\node[
    draw,
    rounded corners=6pt,
    fill=gray!30,
    inner sep=2pt,
    right=0.5em of 12
] {
    \colorbox{gray!30}{first + second constituent}
};
\end{tikzpicture}

\vspace{-0.25em}

\begin{table}[!ht]
\begin{center}
\setlength{\extrarowheight}{1pt}
\begin{tabularx}{\columnwidth}{|X|X|X|}
\hline
 & item & type \\
\hline
 coverage & $NA$ & $NA$ \\
\hline
 regularity & $NA$ & $NA$ \\
\hline
\end{tabularx}
\end{center}
\end{table}

\vspace{-1.35em}

\noindent
Support level: \textit{NA}

\vspace{0.3em}

\noindent
Examples: Erbschaft{\zeroLinker}steuer, Herz{\zeroLinker}schlagen, Schaden{\zeroLinker}ersatz

\noindent
Counterexamples: Schaden\textbf{s}ersatz

\vspace{0.3em}

\noindent
Sources: \citet[11]{nuebling_szczepaniak_2008}, \citet[98]{ortner_mueller-bollhagen_1991}, \citet[20]{schluecker_2023}, \citet[177]{wegener_2005}

\vspace{0.35em}
\hrule




% 39
\vspace{0.4em}

\begin{center}
\textbf{Constraint (39)}
\end{center}

\noindent
Compounds whose first constituent designates a person and whose second constituent is one of 'Mann', 'Frau', 'Leute', 'Tochter', 'Gattin', or 'Witwe' tend to insert tend to insert -s-.
\vspace{0.5em}

\begin{tikzpicture}
\node[
    draw,
    rounded corners=6pt,
    fill=gray!30,
    inner sep=2pt
] (11) {
    \colorbox{gray!30}{P2L}
};
\node[
    draw,
    rounded corners=6pt,
    fill=gray!30,
    inner sep=2pt,
    right=0.5em of 11
] (12) {
    \colorbox{gray!30}{SL}
};
\node[
    draw,
    rounded corners=6pt,
    fill=gray!30,
    inner sep=2pt,
    right=0.5em of 12
] {
    \colorbox{gray!30}{first + second constituent}
};
\end{tikzpicture}

\vspace{-0.25em}

\begin{table}[!ht]
\begin{center}
\setlength{\extrarowheight}{1pt}
\begin{tabularx}{\columnwidth}{|X|X|X|}
\hline
 & item & type \\
\hline
 coverage & $NA$ & $NA$ \\
\hline
 regularity & $NA$ & $NA$ \\
\hline
\end{tabularx}
\end{center}
\end{table}

\vspace{-1.35em}

\noindent
Support level: \textit{NA}

\vspace{0.3em}

\noindent
Examples: Reiter\textbf{s}mann, Lehrer\textbf{s}tochter, Bäcker\textbf{s}leute

\noindent
Counterexamples: 

\vspace{0.3em}

\noindent
Sources: \citet[11]{nuebling_szczepaniak_2008}, \citet[56]{ortner_mueller-bollhagen_1991}

\vspace{0.35em}
\hrule




% 40
\vspace{0.4em}

\begin{center}
\textbf{Constraint (40)}
\end{center}

\noindent
The probability of -s- grows irregularly with the morphological complexity of the first constituent (any form of derivation).
\vspace{0.5em}

\begin{tikzpicture}
\node[
    draw,
    rounded corners=6pt,
    fill=gray!30,
    inner sep=2pt
] (11) {
    \colorbox{gray!30}{P2L}
};
\node[
    draw,
    rounded corners=6pt,
    fill=gray!30,
    inner sep=2pt,
    right=0.5em of 11
] (12) {
    \colorbox{gray!30}{D}
};
\node[
    draw,
    rounded corners=6pt,
    fill=gray!30,
    inner sep=2pt,
    right=0.5em of 12
] {
    \colorbox{gray!30}{first constituent}
};
\end{tikzpicture}

\vspace{-0.25em}

\begin{table}[!ht]
\begin{center}
\setlength{\extrarowheight}{1pt}
\begin{tabularx}{\columnwidth}{|X|X|X|}
\hline
 & item & type \\
\hline
 coverage & $1.0$ & $1.0$ \\
\hline
 regularity & $0.735$ & $0.766$ \\
\hline
\end{tabularx}
\end{center}
\end{table}

\vspace{-1.35em}

\noindent
Support level: \textit{supported}

\vspace{0.3em}

\noindent
Examples: Gesundheit\textbf{s}amt, Gesellschaft\textbf{s}politik, Anspruch\textbf{s}haltung, Eintrag\textbf{s}frist

\noindent
Counterexamples: Ort\textbf{s}amt, Krieg\textbf{s}ende, Anruf{\zeroLinker}beantworter, Minderheit\textbf{en}recht

\vspace{0.3em}

\noindent
Sources: \citet[572]{fuhrhop_kuerschner_2015}, \citet[201]{kopf_2017}, \citet[106, 111, 112]{kuerschner_2005}, \citet[19, 22]{kuerschner_2010}

\vspace{0.35em}
\hrule




% 41
\vspace{0.4em}

\begin{center}
\textbf{Constraint (41)}
\end{center}

\noindent
The probability of -s- grows irregularly with the phonological complexity of the first constituent. By phonologically complex words are understood words with polysyllable non-trochaic form, words with unstressed prefixes, words with stressed or semi-stressed suffixes etc.
\vspace{0.5em}

\begin{tikzpicture}
\node[
    draw,
    rounded corners=6pt,
    fill=gray!30,
    inner sep=2pt
] (11) {
    \colorbox{gray!30}{P2L}
};
\node[
    draw,
    rounded corners=6pt,
    fill=gray!30,
    inner sep=2pt,
    right=0.5em of 11
] (12) {
    \colorbox{gray!30}{P}
};
\node[
    draw,
    rounded corners=6pt,
    fill=gray!30,
    inner sep=2pt,
    right=0.5em of 12
] {
    \colorbox{gray!30}{first constituent}
};
\end{tikzpicture}

\vspace{-0.25em}

\begin{table}[!ht]
\begin{center}
\setlength{\extrarowheight}{1pt}
\begin{tabularx}{\columnwidth}{|X|X|X|}
\hline
 & item & type \\
\hline
 coverage & $1.0$ & $1.0$ \\
\hline
 regularity & $0.723$ & $0.761$ \\
\hline
\end{tabularx}
\end{center}
\end{table}

\vspace{-1.35em}

\noindent
Support level: \textit{supported}

\vspace{0.3em}

\noindent
Examples: Anruf{\zeroLinker}beantworter, Beruf\textbf{s}erfahrung, Religion\textbf{s}unterricht

\noindent
Counterexamples: Schlaf{\zeroLinker}plan, Herbst{\zeroLinker}anfang, Bestand{\zeroLinker}teil

\vspace{0.3em}

\noindent
Sources: \citet[572]{fuhrhop_kuerschner_2015}, \citet[10, 21, 22]{nuebling_szczepaniak_2008}, \citet[78]{nuebling_szczepaniak_2013}

\vspace{0.35em}
\hrule




% 42
\vspace{0.4em}

\begin{center}
\textbf{Constraint (42)}
\end{center}

\noindent
The probability of -s- in compounds with simplex first constituents tends to decrease with the increasing sonority of the final segment of the first constituent. -s- is thus more frequent after plosives, infrequent after nasals and liquids, and it never occurs after a full vowel.
\vspace{0.5em}

\begin{tikzpicture}
\node[
    draw,
    rounded corners=6pt,
    fill=gray!30,
    inner sep=2pt
] (11) {
    \colorbox{gray!30}{P2L}
};
\node[
    draw,
    rounded corners=6pt,
    fill=gray!30,
    inner sep=2pt,
    right=0.5em of 11
] (12) {
    \colorbox{gray!30}{P}
};
\node[
    draw,
    rounded corners=6pt,
    fill=gray!30,
    inner sep=2pt,
    right=0.5em of 12
] {
    \colorbox{gray!30}{first constituent}
};
\end{tikzpicture}

\vspace{-0.25em}

\begin{table}[!ht]
\begin{center}
\setlength{\extrarowheight}{1pt}
\begin{tabularx}{\columnwidth}{|X|X|X|}
\hline
 & item & type \\
\hline
 coverage & $0.431$ & $0.498$ \\
\hline
 regularity & $0.796$ & $0.717$ \\
\hline
\end{tabularx}
\end{center}
\end{table}

\vspace{-1.35em}

\noindent
Support level: \textit{supported}

\vspace{0.3em}

\noindent
Examples: Ort\textbf{s}tarif, Glück\textbf{s}rad, Himmel{\zeroLinker}reich, Uhu{\zeroLinker}paar

\noindent
Counterexamples: Arbeit{\zeroLinker}geber, Dieb{\zeroLinker}stahl, Himmel\textbf{s}bahn, Uhu\textbf{s}nest

\vspace{0.3em}

\noindent
Sources: \citet[572]{fuhrhop_kuerschner_2015}, \citet[110]{kuerschner_2005}, \citet[7, 17]{nuebling_szczepaniak_2008}, \citet[9]{schaefer_pankratz_2018}, \citet[434]{wegener_2003}, \citet[180, 181, 182]{wegener_2005}

\vspace{0.35em}
\hrule




% 43
\vspace{0.4em}

\begin{center}
\textbf{Constraint (43)}
\end{center}

\noindent
Almost all simplex nouns that constitute first constituents that attach -s- belong to a small fixed group of masculine and neuter nouns.
\vspace{0.5em}

\begin{tikzpicture}
\node[
    draw,
    rounded corners=6pt,
    fill=gray!30,
    inner sep=2pt
] (11) {
    \colorbox{gray!30}{L2P}
};
\node[
    draw,
    rounded corners=6pt,
    fill=gray!30,
    inner sep=2pt,
    right=0.5em of 11
] (12) {
    \colorbox{gray!30}{M}
};
\node[
    draw,
    rounded corners=6pt,
    fill=gray!30,
    inner sep=2pt,
    right=0.5em of 12
] {
    \colorbox{gray!30}{first constituent + linker}
};
\end{tikzpicture}

\vspace{-0.25em}

\begin{table}[!ht]
\begin{center}
\setlength{\extrarowheight}{1pt}
\begin{tabularx}{\columnwidth}{|X|X|X|}
\hline
 & item & type \\
\hline
 coverage & $0.05$ & $0.032$ \\
\hline
 regularity & $0.988$ & $0.915$ \\
\hline
\end{tabularx}
\end{center}
\end{table}

\vspace{-1.35em}

\noindent
Support level: \textit{supported}

\vspace{0.3em}

\noindent
Examples: Zwang\textbf{s}jacke, Ort\textbf{s}angabe, König\textbf{s}haus

\noindent
Counterexamples: Arbeit\textbf{s}amt, Heirat\textbf{s}antrag, Liebe\textbf{s}brief

\vspace{0.3em}

\noindent
Sources: \citet[23]{kuerschner_2010}, \citet[79]{nuebling_szczepaniak_2013}, \citet[8]{schaefer_pankratz_2018}

\vspace{0.35em}
\hrule




% 44
\vspace{0.4em}

\begin{center}
\textbf{Constraint (44)}
\end{center}

\noindent
Many simplex nouns that constitute first constituents that attach -s- have a high token frequency.
\vspace{0.5em}

\begin{tikzpicture}
\node[
    draw,
    rounded corners=6pt,
    fill=gray!30,
    inner sep=2pt
] (11) {
    \colorbox{gray!30}{L2P}
};
\node[
    draw,
    rounded corners=6pt,
    fill=gray!30,
    inner sep=2pt,
    right=0.5em of 11
] (12) {
    \colorbox{gray!30}{L}
};
\node[
    draw,
    rounded corners=6pt,
    fill=gray!30,
    inner sep=2pt,
    right=0.5em of 12
] {
    \colorbox{gray!30}{first constituent + linker}
};
\end{tikzpicture}

\vspace{-0.25em}

\begin{table}[!ht]
\begin{center}
\setlength{\extrarowheight}{1pt}
\begin{tabularx}{\columnwidth}{|X|X|X|}
\hline
 & item & type \\
\hline
 coverage & $0.05$ & $0.032$ \\
\hline
 regularity & $0.649$ & $0.914$ \\
\hline
\end{tabularx}
\end{center}
\end{table}

\vspace{-1.35em}

\noindent
Support level: \textit{supported}

\vspace{0.3em}

\noindent
Examples: Volk\textbf{s}brauch, Amt\textbf{s}anwalt, Staat\textbf{s}amt

\noindent
Counterexamples: Fasching\textbf{s}foto, Gral\textbf{s}sage, Konvent\textbf{s}sitzung

\vspace{0.3em}

\noindent
Sources: \citet[25]{kuerschner_2010}

\vspace{0.35em}
\hrule




% 45
\vspace{0.4em}

\begin{center}
\textbf{Constraint (45)}
\end{center}

\noindent
All feminine nouns that constitute first constituents that attach -s- are morphologically complex.
\vspace{0.5em}

\begin{tikzpicture}
\node[
    draw,
    rounded corners=6pt,
    fill=gray!30,
    inner sep=2pt
] (11) {
    \colorbox{gray!30}{L2P}
};
\node[
    draw,
    rounded corners=6pt,
    fill=gray!30,
    inner sep=2pt,
    right=0.5em of 11
] (12) {
    \colorbox{gray!30}{MD}
};
\node[
    draw,
    rounded corners=6pt,
    fill=gray!30,
    inner sep=2pt,
    right=0.5em of 12
] {
    \colorbox{gray!30}{first constituent + linker}
};
\end{tikzpicture}

\vspace{-0.25em}

\begin{table}[!ht]
\begin{center}
\setlength{\extrarowheight}{1pt}
\begin{tabularx}{\columnwidth}{|X|X|X|}
\hline
 & item & type \\
\hline
 coverage & $0.163$ & $0.159$ \\
\hline
 regularity & $0.996$ & $0.983$ \\
\hline
\end{tabularx}
\end{center}
\end{table}

\vspace{-1.35em}

\noindent
Support level: \textit{supported}

\vspace{0.3em}

\noindent
Examples: Gesundheit\textbf{s}amt, Gesellschaft\textbf{s}politik, Liebe\textbf{s}brief

\noindent
Counterexamples: Arbeit\textbf{s}tag, Heirat\textbf{s}antrag

\vspace{0.3em}

\noindent
Sources: \citet[23, 24]{kuerschner_2010}

\vspace{0.35em}
\hrule




% 46
\vspace{0.4em}

\begin{center}
\textbf{Constraint (46)}
\end{center}

\noindent
Almost all feminine nouns that constitute first constituents that attach -s- are polysyllabic.
\vspace{0.5em}

\begin{tikzpicture}
\node[
    draw,
    rounded corners=6pt,
    fill=gray!30,
    inner sep=2pt
] (11) {
    \colorbox{gray!30}{L2P}
};
\node[
    draw,
    rounded corners=6pt,
    fill=gray!30,
    inner sep=2pt,
    right=0.5em of 11
] (12) {
    \colorbox{gray!30}{MP}
};
\node[
    draw,
    rounded corners=6pt,
    fill=gray!30,
    inner sep=2pt,
    right=0.5em of 12
] {
    \colorbox{gray!30}{first constituent + linker}
};
\end{tikzpicture}

\vspace{-0.25em}

\begin{table}[!ht]
\begin{center}
\setlength{\extrarowheight}{1pt}
\begin{tabularx}{\columnwidth}{|X|X|X|}
\hline
 & item & type \\
\hline
 coverage & $0.163$ & $0.159$ \\
\hline
 regularity & $0.999$ & $1.0$ \\
\hline
\end{tabularx}
\end{center}
\end{table}

\vspace{-1.35em}

\noindent
Support level: \textit{supported}

\vspace{0.3em}

\noindent
Examples: Gesundheit\textbf{s}amt, Gesellschaft\textbf{s}politik, Heirat\textbf{s}antrag, Liebe\textbf{s}brief

\noindent
Counterexamples: 

\vspace{0.3em}

\noindent
Sources: \citet[73]{ortner_mueller-bollhagen_1991}

\vspace{0.35em}
\hrule




% 47
\vspace{0.4em}

\begin{center}
\textbf{Constraint (47)}
\end{center}

\noindent
All nouns that constitute first constituents that attach the -n- allomorph of the -(e)n- linker end in schwa.
\vspace{0.5em}

\begin{tikzpicture}
\node[
    draw,
    rounded corners=6pt,
    fill=gray!30,
    inner sep=2pt
] (11) {
    \colorbox{gray!30}{L2P}
};
\node[
    draw,
    rounded corners=6pt,
    fill=gray!30,
    inner sep=2pt,
    right=0.5em of 11
] (12) {
    \colorbox{gray!30}{P}
};
\node[
    draw,
    rounded corners=6pt,
    fill=gray!30,
    inner sep=2pt,
    right=0.5em of 12
] {
    \colorbox{gray!30}{first constituent + linker}
};
\end{tikzpicture}

\vspace{-0.25em}

\begin{table}[!ht]
\begin{center}
\setlength{\extrarowheight}{1pt}
\begin{tabularx}{\columnwidth}{|X|X|X|}
\hline
 & item & type \\
\hline
 coverage & $0.116$ & $0.109$ \\
\hline
 regularity & $0.959$ & $0.985$ \\
\hline
\end{tabularx}
\end{center}
\end{table}

\vspace{-1.35em}

\noindent
Support level: \textit{partially supported}

\vspace{0.3em}

\noindent
Examples: Biene\textbf{n}zucht, Suppe\textbf{n}schüssel, Auge\textbf{n}lied

\noindent
Counterexamples: Oper\textbf{n}haus, Elster\textbf{n}nest

\vspace{0.3em}

\noindent
Sources: \citet[18]{kuerschner_2010}, \citet[84]{nuebling_szczepaniak_2013}, \citet[179]{wegener_2005}

\vspace{0.35em}
\hrule




% 48
\vspace{0.4em}

\begin{center}
\textbf{Constraint (48)}
\end{center}

\noindent
All nouns that constitute first constituents that attach -(e)n- belong to weak nouns.
\vspace{0.5em}

\begin{tikzpicture}
\node[
    draw,
    rounded corners=6pt,
    fill=gray!30,
    inner sep=2pt
] (11) {
    \colorbox{gray!30}{L2P}
};
\node[
    draw,
    rounded corners=6pt,
    fill=gray!30,
    inner sep=2pt,
    right=0.5em of 11
] (12) {
    \colorbox{gray!30}{M}
};
\node[
    draw,
    rounded corners=6pt,
    fill=gray!30,
    inner sep=2pt,
    right=0.5em of 12
] {
    \colorbox{gray!30}{first constituent + linker}
};
\end{tikzpicture}

\vspace{-0.25em}

\begin{table}[!ht]
\begin{center}
\setlength{\extrarowheight}{1pt}
\begin{tabularx}{\columnwidth}{|X|X|X|}
\hline
 & item & type \\
\hline
 coverage & $0.143$ & $0.127$ \\
\hline
 regularity & $0.971$ & $0.986$ \\
\hline
\end{tabularx}
\end{center}
\end{table}

\vspace{-1.35em}

\noindent
Support level: \textit{partially supported}

\vspace{0.3em}

\noindent
Examples: Blume\textbf{n}topf, Katze\textbf{n}fell, Frau\textbf{en}hand

\noindent
Counterexamples: Hahn\textbf{en}kamm, Mond\textbf{en}schein, Instrument\textbf{en}bau

\vspace{0.3em}

\noindent
Sources: \citet[187]{kopf_2017}, \citet[118, 119]{kuerschner_2005}, \citet[9]{kuerschner_2010}, \citet[3]{nuebling_szczepaniak_2008}, \citet[79, 84]{nuebling_szczepaniak_2013}

\vspace{0.35em}
\hrule




% 49
\vspace{0.4em}

\begin{center}
\textbf{Constraint (49)}
\end{center}

\noindent
All masculine nouns that do not build the plural form with -(e)n that constitute first constituents that attach -(e)n- are monosyllabic designations of male persons, animals, astronomic objects, months.
\vspace{0.5em}

\begin{tikzpicture}
\node[
    draw,
    rounded corners=6pt,
    fill=gray!30,
    inner sep=2pt
] (11) {
    \colorbox{gray!30}{L2P}
};
\node[
    draw,
    rounded corners=6pt,
    fill=gray!30,
    inner sep=2pt,
    right=0.5em of 11
] (12) {
    \colorbox{gray!30}{MPS}
};
\node[
    draw,
    rounded corners=6pt,
    fill=gray!30,
    inner sep=2pt,
    right=0.5em of 12
] {
    \colorbox{gray!30}{first constituent + linker}
};
\end{tikzpicture}

\vspace{-0.25em}

\begin{table}[!ht]
\begin{center}
\setlength{\extrarowheight}{1pt}
\begin{tabularx}{\columnwidth}{|X|X|X|}
\hline
 & item & type \\
\hline
 coverage & $NA$ & $NA$ \\
\hline
 regularity & $NA$ & $NA$ \\
\hline
\end{tabularx}
\end{center}
\end{table}

\vspace{-1.35em}

\noindent
Support level: \textit{NA}

\vspace{0.3em}

\noindent
Examples: Hahn\textbf{en}kamm, Mond\textbf{en}schein, Stern\textbf{en}himmel, Mai\textbf{en}nacht

\noindent
Counterexamples: Mond{\zeroLinker}aufgang, Stern{\zeroLinker}bild, Mai{\zeroLinker}baum

\vspace{0.3em}

\noindent
Sources: \citet[79, 80]{nuebling_szczepaniak_2013}, \citet[77, 78]{ortner_mueller-bollhagen_1991}

\vspace{0.35em}
\hrule




% 50
\vspace{0.4em}

\begin{center}
\textbf{Constraint (50)}
\end{center}

\noindent
All neuter nouns that do not build the plural form with -(e)n that constitute first constituents that attach -(e)n- are polysyllabic foreign nouns that have a stressed final syllable that exhibit a plural meaning within the compound.
\vspace{0.5em}

\begin{tikzpicture}
\node[
    draw,
    rounded corners=6pt,
    fill=gray!30,
    inner sep=2pt
] (11) {
    \colorbox{gray!30}{L2P}
};
\node[
    draw,
    rounded corners=6pt,
    fill=gray!30,
    inner sep=2pt,
    right=0.5em of 11
] (12) {
    \colorbox{gray!30}{MPLS}
};
\node[
    draw,
    rounded corners=6pt,
    fill=gray!30,
    inner sep=2pt,
    right=0.5em of 12
] {
    \colorbox{gray!30}{first constituent + linker}
};
\end{tikzpicture}

\vspace{-0.25em}

\begin{table}[!ht]
\begin{center}
\setlength{\extrarowheight}{1pt}
\begin{tabularx}{\columnwidth}{|X|X|X|}
\hline
 & item & type \\
\hline
 coverage & $NA$ & $NA$ \\
\hline
 regularity & $NA$ & $NA$ \\
\hline
\end{tabularx}
\end{center}
\end{table}

\vspace{-1.35em}

\noindent
Support level: \textit{NA}

\vspace{0.3em}

\noindent
Examples: Instrument\textbf{en}bau, Dokument\textbf{en}sammlung

\noindent
Counterexamples: Zertifikat{\zeroLinker}system

\vspace{0.3em}

\noindent
Sources: \citet[78, 79]{ortner_mueller-bollhagen_1991}

\vspace{0.35em}
\hrule




% 51
\vspace{0.4em}

\begin{center}
\textbf{Constraint (51)}
\end{center}

\noindent
All nouns that constitute first constituents that attach -e- build the plural form with -e.
\vspace{0.5em}

\begin{tikzpicture}
\node[
    draw,
    rounded corners=6pt,
    fill=gray!30,
    inner sep=2pt
] (11) {
    \colorbox{gray!30}{L2P}
};
\node[
    draw,
    rounded corners=6pt,
    fill=gray!30,
    inner sep=2pt,
    right=0.5em of 11
] (12) {
    \colorbox{gray!30}{M}
};
\node[
    draw,
    rounded corners=6pt,
    fill=gray!30,
    inner sep=2pt,
    right=0.5em of 12
] {
    \colorbox{gray!30}{first constituent + linker}
};
\end{tikzpicture}

\vspace{-0.25em}

\begin{table}[!ht]
\begin{center}
\setlength{\extrarowheight}{1pt}
\begin{tabularx}{\columnwidth}{|X|X|X|}
\hline
 & item & type \\
\hline
 coverage & $0.023$ & $0.008$ \\
\hline
 regularity & $0.982$ & $0.994$ \\
\hline
\end{tabularx}
\end{center}
\end{table}

\vspace{-1.35em}

\noindent
Support level: \textit{supported}

\vspace{0.3em}

\noindent
Examples: Tag\textbf{e}buch, Punkt\textbf{e}stand, Hund\textbf{e}leine

\noindent
Counterexamples: Herz\textbf{e}leid, Maus\textbf{e}zahn

\vspace{0.3em}

\noindent
Sources: \citet[9, 18]{kuerschner_2010}, \citet[5]{nuebling_szczepaniak_2008}, \citet[70, 81]{nuebling_szczepaniak_2013}, \citet[102]{ortner_mueller-bollhagen_1991}

\vspace{0.35em}
\hrule




% 52
\vspace{0.4em}

\begin{center}
\textbf{Constraint (52)}
\end{center}

\noindent
All nouns that constitute first constituents that attach -e- have a stressed last syllable.
\vspace{0.5em}

\begin{tikzpicture}
\node[
    draw,
    rounded corners=6pt,
    fill=gray!30,
    inner sep=2pt
] (11) {
    \colorbox{gray!30}{L2P}
};
\node[
    draw,
    rounded corners=6pt,
    fill=gray!30,
    inner sep=2pt,
    right=0.5em of 11
] (12) {
    \colorbox{gray!30}{P}
};
\node[
    draw,
    rounded corners=6pt,
    fill=gray!30,
    inner sep=2pt,
    right=0.5em of 12
] {
    \colorbox{gray!30}{first constituent + linker}
};
\end{tikzpicture}

\vspace{-0.25em}

\begin{table}[!ht]
\begin{center}
\setlength{\extrarowheight}{1pt}
\begin{tabularx}{\columnwidth}{|X|X|X|}
\hline
 & item & type \\
\hline
 coverage & $0.023$ & $0.008$ \\
\hline
 regularity & $0.973$ & $0.992$ \\
\hline
\end{tabularx}
\end{center}
\end{table}

\vspace{-1.35em}

\noindent
Support level: \textit{partially supported}

\vspace{0.3em}

\noindent
Examples: Tag\textbf{e}buch, Punkt\textbf{e}stand, Hund\textbf{e}leine

\noindent
Counterexamples: Anhalt\textbf{e}stelle

\vspace{0.3em}

\noindent
Sources: \citet[9, 18]{kuerschner_2010}

\vspace{0.35em}
\hrule




% 53
\vspace{0.4em}

\begin{center}
\textbf{Constraint (53)}
\end{center}

\noindent
Most nouns that constitute first constituents that attach -e- are simplex.
\vspace{0.5em}

\begin{tikzpicture}
\node[
    draw,
    rounded corners=6pt,
    fill=gray!30,
    inner sep=2pt
] (11) {
    \colorbox{gray!30}{L2P}
};
\node[
    draw,
    rounded corners=6pt,
    fill=gray!30,
    inner sep=2pt,
    right=0.5em of 11
] (12) {
    \colorbox{gray!30}{D}
};
\node[
    draw,
    rounded corners=6pt,
    fill=gray!30,
    inner sep=2pt,
    right=0.5em of 12
] {
    \colorbox{gray!30}{first constituent + linker}
};
\end{tikzpicture}

\vspace{-0.25em}

\begin{table}[!ht]
\begin{center}
\setlength{\extrarowheight}{1pt}
\begin{tabularx}{\columnwidth}{|X|X|X|}
\hline
 & item & type \\
\hline
 coverage & $0.023$ & $0.008$ \\
\hline
 regularity & $0.73$ & $0.783$ \\
\hline
\end{tabularx}
\end{center}
\end{table}

\vspace{-1.35em}

\noindent
Support level: \textit{supported}

\vspace{0.3em}

\noindent
Examples: Tag\textbf{e}buch, Punkt\textbf{e}stand, Hund\textbf{e}leine

\noindent
Counterexamples: Anhalt\textbf{e}stelle

\vspace{0.3em}

\noindent
Sources: \citet[9, 18]{kuerschner_2010}, \citet[102]{ortner_mueller-bollhagen_1991}

\vspace{0.35em}
\hrule




% 54
\vspace{0.4em}

\begin{center}
\textbf{Constraint (54)}
\end{center}

\noindent
All nouns that constitute first constituents that attach -e- are native (none are loanwords).
\vspace{0.5em}

\begin{tikzpicture}
\node[
    draw,
    rounded corners=6pt,
    fill=gray!30,
    inner sep=2pt
] (11) {
    \colorbox{gray!30}{L2P}
};
\node[
    draw,
    rounded corners=6pt,
    fill=gray!30,
    inner sep=2pt,
    right=0.5em of 11
] (12) {
    \colorbox{gray!30}{L}
};
\node[
    draw,
    rounded corners=6pt,
    fill=gray!30,
    inner sep=2pt,
    right=0.5em of 12
] {
    \colorbox{gray!30}{first constituent + linker}
};
\end{tikzpicture}

\vspace{-0.25em}

\begin{table}[!ht]
\begin{center}
\setlength{\extrarowheight}{1pt}
\begin{tabularx}{\columnwidth}{|X|X|X|}
\hline
 & item & type \\
\hline
 coverage & $NA$ & $NA$ \\
\hline
 regularity & $NA$ & $NA$ \\
\hline
\end{tabularx}
\end{center}
\end{table}

\vspace{-1.35em}

\noindent
Support level: \textit{NA}

\vspace{0.3em}

\noindent
Examples: Tag\textbf{e}buch, Punkt\textbf{e}stand, Hund\textbf{e}leine

\noindent
Counterexamples: 

\vspace{0.3em}

\noindent
Sources: \citet[102]{ortner_mueller-bollhagen_1991}

\vspace{0.35em}
\hrule




% 55
\vspace{0.4em}

\begin{center}
\textbf{Constraint (55)}
\end{center}

\noindent
All nouns that constitute first constituents that attach -(")er- build the plural form with -(")er.
\vspace{0.5em}

\begin{tikzpicture}
\node[
    draw,
    rounded corners=6pt,
    fill=gray!30,
    inner sep=2pt
] (11) {
    \colorbox{gray!30}{L2P}
};
\node[
    draw,
    rounded corners=6pt,
    fill=gray!30,
    inner sep=2pt,
    right=0.5em of 11
] (12) {
    \colorbox{gray!30}{M}
};
\node[
    draw,
    rounded corners=6pt,
    fill=gray!30,
    inner sep=2pt,
    right=0.5em of 12
] {
    \colorbox{gray!30}{first constituent + linker}
};
\end{tikzpicture}

\vspace{-0.25em}

\begin{table}[!ht]
\begin{center}
\setlength{\extrarowheight}{1pt}
\begin{tabularx}{\columnwidth}{|X|X|X|}
\hline
 & item & type \\
\hline
 coverage & $0.01$ & $0.015$ \\
\hline
 regularity & $0.979$ & $0.999$ \\
\hline
\end{tabularx}
\end{center}
\end{table}

\vspace{-1.35em}

\noindent
Support level: \textit{supported}

\vspace{0.3em}

\noindent
Examples: Buch\textbf{=er}regal, Kraut\textbf{=er}frau, Kind\textbf{=er}jacke

\noindent
Counterexamples: 

\vspace{0.3em}

\noindent
Sources: \citet[229]{eisenberg_2013}, \citet[9, 11, 18]{kuerschner_2010}, \citet[3]{nuebling_szczepaniak_2008}, \citet[80, 81]{nuebling_szczepaniak_2013}, \citet[98]{ortner_mueller-bollhagen_1991}

\vspace{0.35em}
\hrule




% 56
\vspace{0.4em}

\begin{center}
\textbf{Constraint (56)}
\end{center}

\noindent
All nouns that constitute first constituents that attach -(")er- have a stressed last syllable.
\vspace{0.5em}

\begin{tikzpicture}
\node[
    draw,
    rounded corners=6pt,
    fill=gray!30,
    inner sep=2pt
] (11) {
    \colorbox{gray!30}{L2P}
};
\node[
    draw,
    rounded corners=6pt,
    fill=gray!30,
    inner sep=2pt,
    right=0.5em of 11
] (12) {
    \colorbox{gray!30}{P}
};
\node[
    draw,
    rounded corners=6pt,
    fill=gray!30,
    inner sep=2pt,
    right=0.5em of 12
] {
    \colorbox{gray!30}{first constituent + linker}
};
\end{tikzpicture}

\vspace{-0.25em}

\begin{table}[!ht]
\begin{center}
\setlength{\extrarowheight}{1pt}
\begin{tabularx}{\columnwidth}{|X|X|X|}
\hline
 & item & type \\
\hline
 coverage & $0.01$ & $0.015$ \\
\hline
 regularity & $0.938$ & $0.96$ \\
\hline
\end{tabularx}
\end{center}
\end{table}

\vspace{-1.35em}

\noindent
Support level: \textit{partially supported}

\vspace{0.3em}

\noindent
Examples: Buch\textbf{=er}regal, Kraut\textbf{=er}frau, Kind\textbf{=er}jacke

\noindent
Counterexamples: Mitglied\textbf{=er}liste, Vorbild\textbf{=er}sammlung

\vspace{0.3em}

\noindent
Sources: \citet[9, 11]{kuerschner_2010}

\vspace{0.35em}
\hrule




% 57
\vspace{0.4em}

\begin{center}
\textbf{Constraint (57)}
\end{center}

\noindent
Most nouns that constitute first constituents that attach -(")er- are simplex.
\vspace{0.5em}

\begin{tikzpicture}
\node[
    draw,
    rounded corners=6pt,
    fill=gray!30,
    inner sep=2pt
] (11) {
    \colorbox{gray!30}{L2P}
};
\node[
    draw,
    rounded corners=6pt,
    fill=gray!30,
    inner sep=2pt,
    right=0.5em of 11
] (12) {
    \colorbox{gray!30}{D}
};
\node[
    draw,
    rounded corners=6pt,
    fill=gray!30,
    inner sep=2pt,
    right=0.5em of 12
] {
    \colorbox{gray!30}{first constituent + linker}
};
\end{tikzpicture}

\vspace{-0.25em}

\begin{table}[!ht]
\begin{center}
\setlength{\extrarowheight}{1pt}
\begin{tabularx}{\columnwidth}{|X|X|X|}
\hline
 & item & type \\
\hline
 coverage & $0.01$ & $0.015$ \\
\hline
 regularity & $0.875$ & $0.935$ \\
\hline
\end{tabularx}
\end{center}
\end{table}

\vspace{-1.35em}

\noindent
Support level: \textit{supported}

\vspace{0.3em}

\noindent
Examples: Buch\textbf{=er}regal, Kraut\textbf{=er}frau, Kind\textbf{=er}jacke

\noindent
Counterexamples: Mitglied\textbf{=er}liste, Vorbild\textbf{=er}sammlung

\vspace{0.3em}

\noindent
Sources: \citet[98]{ortner_mueller-bollhagen_1991}

\vspace{0.35em}
\hrule




% 58
\vspace{0.4em}

\begin{center}
\textbf{Constraint (58)}
\end{center}

\noindent
All nouns that constitute first constituents that attach -(")er- are native (none are loanwords).
\vspace{0.5em}

\begin{tikzpicture}
\node[
    draw,
    rounded corners=6pt,
    fill=gray!30,
    inner sep=2pt
] (11) {
    \colorbox{gray!30}{L2P}
};
\node[
    draw,
    rounded corners=6pt,
    fill=gray!30,
    inner sep=2pt,
    right=0.5em of 11
] (12) {
    \colorbox{gray!30}{L}
};
\node[
    draw,
    rounded corners=6pt,
    fill=gray!30,
    inner sep=2pt,
    right=0.5em of 12
] {
    \colorbox{gray!30}{first constituent + linker}
};
\end{tikzpicture}

\vspace{-0.25em}

\begin{table}[!ht]
\begin{center}
\setlength{\extrarowheight}{1pt}
\begin{tabularx}{\columnwidth}{|X|X|X|}
\hline
 & item & type \\
\hline
 coverage & $NA$ & $NA$ \\
\hline
 regularity & $NA$ & $NA$ \\
\hline
\end{tabularx}
\end{center}
\end{table}

\vspace{-1.35em}

\noindent
Support level: \textit{NA}

\vspace{0.3em}

\noindent
Examples: Buch\textbf{=er}regal, Kraut\textbf{=er}frau, Kind\textbf{=er}jacke

\noindent
Counterexamples: 

\vspace{0.3em}

\noindent
Sources: \citet[98]{ortner_mueller-bollhagen_1991}

\vspace{0.35em}
\hrule




% 59
\vspace{0.4em}

\begin{center}
\textbf{Constraint (59)}
\end{center}

\noindent
All nouns that constitute first constituents that attach -"e- build the plural form with -e and umlaut.
\vspace{0.5em}

\begin{tikzpicture}
\node[
    draw,
    rounded corners=6pt,
    fill=gray!30,
    inner sep=2pt
] (11) {
    \colorbox{gray!30}{L2P}
};
\node[
    draw,
    rounded corners=6pt,
    fill=gray!30,
    inner sep=2pt,
    right=0.5em of 11
] (12) {
    \colorbox{gray!30}{M}
};
\node[
    draw,
    rounded corners=6pt,
    fill=gray!30,
    inner sep=2pt,
    right=0.5em of 12
] {
    \colorbox{gray!30}{first constituent + linker}
};
\end{tikzpicture}

\vspace{-0.25em}

\begin{table}[!ht]
\begin{center}
\setlength{\extrarowheight}{1pt}
\begin{tabularx}{\columnwidth}{|X|X|X|}
\hline
 & item & type \\
\hline
 coverage & $0.011$ & $0.003$ \\
\hline
 regularity & $1.0$ & $1.0$ \\
\hline
\end{tabularx}
\end{center}
\end{table}

\vspace{-1.35em}

\noindent
Support level: \textit{supported}

\vspace{0.3em}

\noindent
Examples: Hand\textbf{=e}druck, Arzt\textbf{=e}streik, Gast\textbf{=e}buch

\noindent
Counterexamples: 

\vspace{0.3em}

\noindent
Sources: \citet[228]{eisenberg_2013}, \citet[5]{nuebling_szczepaniak_2008}

\vspace{0.35em}
\hrule




% 60
\vspace{0.4em}

\begin{center}
\textbf{Constraint (60)}
\end{center}

\noindent
All nouns that constitute first constituents that attach -es- belong to a fixed row of masculine and neuter nouns and build the genitive form with -(e)s-. -es- is thus isolated.
\vspace{0.5em}

\begin{tikzpicture}
\node[
    draw,
    rounded corners=6pt,
    fill=gray!30,
    inner sep=2pt
] (11) {
    \colorbox{gray!30}{L2P}
};
\node[
    draw,
    rounded corners=6pt,
    fill=gray!30,
    inner sep=2pt,
    right=0.5em of 11
] (12) {
    \colorbox{gray!30}{LM}
};
\node[
    draw,
    rounded corners=6pt,
    fill=gray!30,
    inner sep=2pt,
    right=0.5em of 12
] {
    \colorbox{gray!30}{first constituent + linker}
};
\end{tikzpicture}

\vspace{-0.25em}

\begin{table}[!ht]
\begin{center}
\setlength{\extrarowheight}{1pt}
\begin{tabularx}{\columnwidth}{|X|X|X|}
\hline
 & item & type \\
\hline
 coverage & $NA$ & $NA$ \\
\hline
 regularity & $NA$ & $NA$ \\
\hline
\end{tabularx}
\end{center}
\end{table}

\vspace{-1.35em}

\noindent
Support level: \textit{NA}

\vspace{0.3em}

\noindent
Examples: Bund\textbf{es}tag, Kind\textbf{es}alter, Tag\textbf{es}schau, Land\textbf{es}regierung

\noindent
Counterexamples: 

\vspace{0.3em}

\noindent
Sources: \citet[229]{eisenberg_2013}, \citet[11, 18]{kuerschner_2010}, \citet[3, 7]{nuebling_szczepaniak_2008}, \citet[81]{nuebling_szczepaniak_2013}, \citet[8]{schaefer_pankratz_2018}

\vspace{0.35em}
\hrule




% 61
\vspace{0.4em}

\begin{center}
\textbf{Constraint (61)}
\end{center}

\noindent
All nouns that constitute first constituents that attach -es- are monosyllabic.
\vspace{0.5em}

\begin{tikzpicture}
\node[
    draw,
    rounded corners=6pt,
    fill=gray!30,
    inner sep=2pt
] (11) {
    \colorbox{gray!30}{L2P}
};
\node[
    draw,
    rounded corners=6pt,
    fill=gray!30,
    inner sep=2pt,
    right=0.5em of 11
] (12) {
    \colorbox{gray!30}{P}
};
\node[
    draw,
    rounded corners=6pt,
    fill=gray!30,
    inner sep=2pt,
    right=0.5em of 12
] {
    \colorbox{gray!30}{first constituent + linker}
};
\end{tikzpicture}

\vspace{-0.25em}

\begin{table}[!ht]
\begin{center}
\setlength{\extrarowheight}{1pt}
\begin{tabularx}{\columnwidth}{|X|X|X|}
\hline
 & item & type \\
\hline
 coverage & $0.012$ & $0.011$ \\
\hline
 regularity & $0.914$ & $0.9$ \\
\hline
\end{tabularx}
\end{center}
\end{table}

\vspace{-1.35em}

\noindent
Support level: \textit{partially supported}

\vspace{0.3em}

\noindent
Examples: Bund\textbf{es}tag, Kind\textbf{es}alter, Tag\textbf{es}schau, Land\textbf{es}regierung

\noindent
Counterexamples: Gesetz\textbf{es}änderung, Vorjahr\textbf{es}winner, Bestand\textbf{es}aufnahme

\vspace{0.3em}

\noindent
Sources: \citet[11, 18]{kuerschner_2010}

\vspace{0.35em}
\hrule




% 62
\vspace{0.4em}

\begin{center}
\textbf{Constraint (62)}
\end{center}

\noindent
All nouns that constitute first constituents that attach -(e)ns- belong to a fixed group of a few masculine nouns and a single neuter noun and have different genitive endings. -(e)ns- is thus isolated.
\vspace{0.5em}

\begin{tikzpicture}
\node[
    draw,
    rounded corners=6pt,
    fill=gray!30,
    inner sep=2pt
] (11) {
    \colorbox{gray!30}{L2P}
};
\node[
    draw,
    rounded corners=6pt,
    fill=gray!30,
    inner sep=2pt,
    right=0.5em of 11
] (12) {
    \colorbox{gray!30}{L}
};
\node[
    draw,
    rounded corners=6pt,
    fill=gray!30,
    inner sep=2pt,
    right=0.5em of 12
] {
    \colorbox{gray!30}{first constituent + linker}
};
\end{tikzpicture}

\vspace{-0.25em}

\begin{table}[!ht]
\begin{center}
\setlength{\extrarowheight}{1pt}
\begin{tabularx}{\columnwidth}{|X|X|X|}
\hline
 & item & type \\
\hline
 coverage & $NA$ & $NA$ \\
\hline
 regularity & $NA$ & $NA$ \\
\hline
\end{tabularx}
\end{center}
\end{table}

\vspace{-1.35em}

\noindent
Support level: \textit{NA}

\vspace{0.3em}

\noindent
Examples: Schmerz\textbf{ens}geld, Name\textbf{ens}tag, Herz\textbf{ens}angst

\noindent
Counterexamples: 

\vspace{0.3em}

\noindent
Sources: \citet[82]{nuebling_szczepaniak_2013}, \citet[80, 97]{ortner_mueller-bollhagen_1991}, \citet[9]{schaefer_pankratz_2018}

\vspace{0.35em}
\hrule




% 63
\vspace{0.4em}

\begin{center}
\textbf{Constraint (63)}
\end{center}

\noindent
All nouns that constitute first constituents that attach a zero linker with umlaut build the plural form with a zero ending and umlaut.
\vspace{0.5em}

\begin{tikzpicture}
\node[
    draw,
    rounded corners=6pt,
    fill=gray!30,
    inner sep=2pt
] (11) {
    \colorbox{gray!30}{L2P}
};
\node[
    draw,
    rounded corners=6pt,
    fill=gray!30,
    inner sep=2pt,
    right=0.5em of 11
] (12) {
    \colorbox{gray!30}{M}
};
\node[
    draw,
    rounded corners=6pt,
    fill=gray!30,
    inner sep=2pt,
    right=0.5em of 12
] {
    \colorbox{gray!30}{first constituent + linker}
};
\end{tikzpicture}

\vspace{-0.25em}

\begin{table}[!ht]
\begin{center}
\setlength{\extrarowheight}{1pt}
\begin{tabularx}{\columnwidth}{|X|X|X|}
\hline
 & item & type \\
\hline
 coverage & $0.003$ & $0.0$ \\
\hline
 regularity & $1.0$ & $1.0$ \\
\hline
\end{tabularx}
\end{center}
\end{table}

\vspace{-1.35em}

\noindent
Support level: \textit{supported}

\vspace{0.3em}

\noindent
Examples: Bruder\textbf{=}gemeinde, Mutter\textbf{=}zentrum, Vater\textbf{=}aufbruch

\noindent
Counterexamples: 

\vspace{0.3em}

\noindent
Sources: \citet[5]{nuebling_szczepaniak_2008}

\vspace{0.35em}
\hrule




% 64
\vspace{0.4em}

\begin{center}
\textbf{Constraint (64)}
\end{center}

\noindent
Linkers that are identical to the plural ending of the first constituent in a given compound tend to be associated with a plural meaning of this first constituent within the compound.
\vspace{0.5em}

\begin{tikzpicture}
\node[
    draw,
    rounded corners=6pt,
    fill=gray!30,
    inner sep=2pt
] (11) {
    \colorbox{gray!30}{L2P}
};
\node[
    draw,
    rounded corners=6pt,
    fill=gray!30,
    inner sep=2pt,
    right=0.5em of 11
] (12) {
    \colorbox{gray!30}{MS}
};
\node[
    draw,
    rounded corners=6pt,
    fill=gray!30,
    inner sep=2pt,
    right=0.5em of 12
] {
    \colorbox{gray!30}{first constituent + linker}
};
\end{tikzpicture}

\vspace{-0.25em}

\begin{table}[!ht]
\begin{center}
\setlength{\extrarowheight}{1pt}
\begin{tabularx}{\columnwidth}{|X|X|X|}
\hline
 & item & type \\
\hline
 coverage & $NA$ & $NA$ \\
\hline
 regularity & $NA$ & $NA$ \\
\hline
\end{tabularx}
\end{center}
\end{table}

\vspace{-1.35em}

\noindent
Support level: \textit{NA}

\vspace{0.3em}

\noindent
Examples: Punkt\textbf{e}stand, Buch\textbf{=er}regal, Staat\textbf{en}bund, Burg\textbf{en}land, Mutter\textbf{=}zentrum

\noindent
Counterexamples: Motor\textbf{en}geräusch, Burg\textbf{en}blick, Kind\textbf{=er}jacke, Hund\textbf{e}leine, Pilz{\zeroLinker}sammler, Buch{\zeroLinker}handel, Vogel{\zeroLinker}futter

\vspace{0.3em}

\noindent
Sources: \citet[3]{krott_etal_2007}, \citet[230, 231]{schaefer_2018}, \citet[20]{schluecker_2023}, \citet[173]{wegener_2005}

\vspace{0.35em}
\hrule




